\documentclass[11pt]{article}

\usepackage{amsmath, amsthm, amssymb}
\usepackage{mathtools}
\usepackage{booktabs}
\usepackage{geometry}
\usepackage{hyperref}
\usepackage{microtype}
\usepackage{enumitem}
\usepackage{xcolor}

\geometry{margin=1.2in}

\newtheorem{theorem}{Theorem}
\newtheorem{lemma}[theorem]{Lemma}
\newtheorem{proposition}[theorem]{Proposition}
\newtheorem{corollary}[theorem]{Corollary}
\newtheorem{remark}[theorem]{Remark}
\theoremstyle{definition}
\newtheorem{assumption}{Assumption}
\newtheorem{definition}[theorem]{Definition}

\newcommand{\E}{\mathbb{E}}
\newcommand{\R}{\mathbb{R}}
\newcommand{\Prob}{\mathbb{P}}
\newcommand{\norm}[1]{\left\| #1 \right\|}
\newcommand{\abs}[1]{\left| #1 \right|}
\newcommand{\inner}[2]{\left\langle #1,\, #2 \right\rangle}
\newcommand{\clip}{\operatorname{Clip}}
\newcommand{\phat}{\hat{p}}
\newcommand{\shat}{\hat{\sigma}}
\newcommand{\ahat}{\hat{\alpha}}
\newcommand{\Filt}{\mathcal{F}}

\title{\textbf{Convergence Analysis of $\alpha$-Robust SGD\\[4pt]
\large Adaptive Gradient Clipping via Online Tail Index Estimation}}
\author{}
\date{}

\begin{document}
\maketitle
\tableofcontents
\newpage

% ===========================================================================
\section{Setup and Notation}
\label{sec:setup}
% ===========================================================================

\paragraph{Problem.} We study the non-convex stochastic optimization problem:
\[
  \min_{x \in \R^d} f(x) \;=\; \E_\xi[F(x, \xi)].
\]

We make the following standing assumptions throughout.

\begin{assumption}[$L$-smoothness]\label{ass:smooth}
$f \colon \R^d \to \R$ is differentiable and $L$-smooth: for all $x, y \in \R^d$,
\[
  \norm{\nabla f(x) - \nabla f(y)} \leq L\norm{x - y}.
\]
Equivalently, for all $x, y$:
$f(y) \leq f(x) + \inner{\nabla f(x)}{y - x} + \frac{L}{2}\norm{y - x}^2$.
\end{assumption}

\begin{assumption}[Heavy-tailed unbiased gradient noise]\label{ass:heavy}
The stochastic gradient $g(x, \xi)$ satisfies, with
$\varepsilon_t := g(x_t, \xi_t) - \nabla f(x_t)$ for the iterates $\{x_t\}$
produced by the algorithm:
\begin{enumerate}[label=(\alph*),noitemsep,leftmargin=*]
  \item \textbf{Unbiasedness.} $\E[\varepsilon_t \mid \Filt_t] = 0$, where
    $\Filt_t := \sigma(x_1, \varepsilon_1, \ldots, \varepsilon_{t-1})$ is the
    natural filtration of the iterates (so $x_t \in \Filt_t$).
  \item \textbf{Bounded $p$-th moment.} There exist $p \in (1, 2]$ and
    $\sigma > 0$ with
    $\E\!\left[\norm{\varepsilon_t}^p \,\big|\, \Filt_t\right] \leq \sigma^p$
    for all $t$.
  \item \textbf{Conditional independence.} Given the trajectory
    $\{x_s\}_{s \leq T}$, the noise terms $\{\varepsilon_t\}$ are independent.
\end{enumerate}
When $p < 2$, the conditional variance $\E[\norm{\varepsilon_t}^2 \mid \Filt_t]$
may be infinite---this is the heavy-tailed regime.
\end{assumption}

\begin{assumption}[Bounded below]\label{ass:below}
$f^* := \inf_{x \in \R^d} f(x) > -\infty$.
We write $\Delta_f := f(x_1) - f^*$ for the initial optimality gap.
\end{assumption}

\begin{assumption}[Bounded gradient signal]\label{ass:signal}
There exists $M > 0$ such that $\norm{\nabla f(x_t)} \leq M$ for all
iterates $x_t$ produced by the algorithm. This is mild: it is implied by
bounded iterates $\norm{x_t} \leq R$ together with $L$-smoothness, or by
compact sublevel sets of $f$.
\end{assumption}

\begin{assumption}[Uniformly regularly varying noise tails]\label{ass:tail}
There exist $\alpha > 1$, $x_0 > 0$, and a slowly varying function $L$
(i.e., $L(tx)/L(x) \to 1$ for all $t > 0$) such that, \emph{uniformly}
over all iterates $x_t$ produced by the algorithm:
\[
  \Prob(\norm{\varepsilon_t} > x \mid x_t) \leq C_{\mathrm{tail}}\, x^{-\alpha} L(x),
  \qquad x \geq x_0,
\]
and there exists $c_{\mathrm{tail}} > 0$ such that:
\[
  \Prob(\norm{\varepsilon_t} > x \mid x_t) \geq c_{\mathrm{tail}}\, x^{-\alpha} L(x),
  \qquad x \geq x_0,
\]
for constants $C_{\mathrm{tail}} \geq c_{\mathrm{tail}} > 0$ that do not depend
on $x_t$. The tail index $\alpha$ and the moment exponent $p$ from
Assumption~\ref{ass:heavy} satisfy $p < \alpha$ (since the $p$-th moment
is finite iff $p < \alpha$).

This is weaker than requiring a common marginal distribution: the noise
may depend on $x_t$ as long as the tail behavior is uniformly controlled.
In our synthetic experiments (Phase~1), the noise is drawn i.i.d.\ from an
exact Pareto distribution independent of $x_t$, so this assumption holds
with $C_{\mathrm{tail}} = c_{\mathrm{tail}} = 1$.

This assumption is needed \emph{only} for the Hill estimator analysis
(Sections~\ref{sec:hill} and~\ref{sec:growing}). The fixed-schedule convergence
theorem (Theorem~\ref{thm:fixed}) requires only
Assumptions~\ref{ass:smooth}--\ref{ass:signal}.
\end{assumption}

\paragraph{Algorithm ($\alpha$-Robust SGD; lagged-window contract).}
Given a time horizon $T$, the algorithm produces iterates:
\[
  x_{t+1} = x_t - \eta_t \cdot \clip(g_t,\, \tau_t), \qquad t = 1, \ldots, T,
\]
where:
\begin{itemize}[noitemsep]
  \item $\clip(g, \tau) := g \cdot \min(1,\, \tau/\norm{g})$ is the norm-clipping
    operator, satisfying $\norm{\clip(g, \tau)} \leq \tau$ for all $g$.
  \item $\eta_t > 0$ is the learning rate (constant under \emph{fixed} schedule;
    $\eta_t = \eta_0/\sqrt{t}$ under \emph{growing} schedule).
  \item $\tau_t > 0$ is the clipping threshold; in the \textbf{adaptive variant},
    $\tau_t$ is computed from the \emph{lagged} rolling window of past gradient
    norms $\{\norm{g_s}\}_{s = \max(1, t-W)}^{t-1}$ (and the Hill estimate
    $\phat_t$ on the same window). The window is updated only \emph{after}
    $\tau_t$ is used at step~$t$, so $\tau_t$ is $\Filt_t$-measurable. See
    \texttt{optimizers/alpha\_robust\_sgd.py:step} and
    \texttt{optimizers/torch/alpha\_robust\_sgd.py:step}.
\end{itemize}

\begin{remark}[Why lag the window]\label{rem:lag}
If $\norm{g_t}$ were included in the window used to compute $\tau_t$, then
$\tau_t$ would be a measurable function of $\varepsilon_t$ and so $\tau_t \notin
\Filt_t$. The clipping inequality $\norm{b_t}\leq B(\tau)$ in
Lemma~\ref{lem:bias} requires $\tau$ to be $\Filt_t$-measurable to apply
$\E[\,\cdot\mid\Filt_t]$. The lagged-window contract makes the conditional
expectation $\E[g_t^c \mid \Filt_t]$ a deterministic function of the
$\Filt_t$-measurable $(x_t,\tau_t)$, recovering Lemmas~\ref{lem:bias}--\ref{lem:variance}
verbatim with $\tau \rightsquigarrow \tau_t$.
\end{remark}

\paragraph{Target.} We will show that $\alpha$-Robust SGD achieves:
\[
  \frac{1}{T}\sum_{t=1}^{T}\E\!\left[\norm{\nabla f(x_t)}^2\right]
  \;\leq\;
  O\!\left(T^{-\frac{2(p-1)}{3p-2}}\right),
\]
which matches the minimax-optimal rate established in Zhang et al.~(2020)
for non-convex optimization under heavy-tailed noise with bounded $p$-th moment.

\begin{remark}[Rate exponent]
For $p = 2$ (bounded variance), the rate reduces to $O(T^{-1/2})$, recovering
the classical non-convex SGD rate. As $p \to 1^+$, the rate degrades to $O(1)$
(no convergence), which is information-theoretically necessary since the first
moment barely exists.
\end{remark}


% ===========================================================================
\section{Clipped Gradient: Bias and Variance Bounds}
\label{sec:clipping}
% ===========================================================================

The next two lemmas are the standard clipping bounds we use in the descent step.

\begin{lemma}[Pointwise clipping inequality]\label{lem:pointwise}
For any $a \geq 0$, $\tau > 0$, and $p \in (1, 2]$:
\begin{enumerate}[label=(\alph*)]
  \item $a \cdot \mathbf{1}(a > \tau) \;\leq\; \dfrac{a^p}{\tau^{p-1}}$.
  \item $\min(a, \tau)^2 \;\leq\; \tau^{2-p}\, a^p$.
\end{enumerate}
\end{lemma}

\begin{proof}
\textbf{(a)} On the event $\{a > \tau\}$: since $1 - p < 0$ and $a > \tau > 0$,
we have $a^{1-p} < \tau^{1-p}$. Multiplying both sides by $a^p > 0$:
\[
  a = a^p \cdot a^{1-p} < a^p \cdot \tau^{1-p} = \frac{a^p}{\tau^{p-1}}.
\]
On $\{a \leq \tau\}$: the left side is $0 \leq a^p/\tau^{p-1}$.

\textbf{(b)} Consider two cases:
\begin{itemize}
  \item If $a \leq \tau$: $\min(a,\tau)^2 = a^2 = a^p \cdot a^{2-p}
    \leq a^p \cdot \tau^{2-p}$, since $a \leq \tau$ and $2 - p \geq 0$.
  \item If $a > \tau$: $\min(a,\tau)^2 = \tau^2 = \tau^{2-p} \cdot \tau^p
    \leq \tau^{2-p} \cdot a^p$, since $a > \tau$.
    \qedhere
\end{itemize}
\end{proof}

\begin{lemma}[Clipping bias bound]\label{lem:bias}
Under Assumptions~\ref{ass:heavy} and~\ref{ass:signal}, let
$g_t = \nabla f(x_t) + \varepsilon_t$ be the stochastic gradient at step $t$,
with $\E[\varepsilon_t \mid x_t] = 0$ and $\E[\norm{\varepsilon_t}^p \mid x_t] \leq \sigma^p$.
Define the clipped gradient $g_t^c = \clip(g_t, \tau)$ and the clipping bias:
\[
  b_t \;:=\; \nabla f(x_t) - \E[g_t^c \mid x_t]
  \;=\; \E\!\left[g_t - g_t^c \mid x_t\right].
\]
Then for any $\tau > M$ (where $M = \sup_t \norm{\nabla f(x_t)}$):
\[
  \norm{b_t} \;\leq\; \frac{\sigma^p}{(\tau - M)^{p-1}}.
\]
In particular, if $\tau \geq 2M$:
$\;\norm{b_t} \leq 2^{p-1}\sigma^p / \tau^{p-1}$.
\end{lemma}

\begin{proof}
The clipping error satisfies:
\[
  \norm{g_t - g_t^c}
  = \norm{g_t}\!\left(1 - \min\!\left(1, \frac{\tau}{\norm{g_t}}\right)\right)
  = \bigl(\norm{g_t} - \tau\bigr)^+.
\]
By Jensen's inequality (applied to the convex function $\norm{\cdot}$):
\[
  \norm{b_t} = \norm{\E[g_t - g_t^c \mid x_t]}
  \leq \E\!\left[\norm{g_t - g_t^c} \mid x_t\right]
  = \E\!\left[\bigl(\norm{g_t} - \tau\bigr)^+ \mid x_t\right].
\]

\textbf{Key step.} On the event $\{\norm{g_t} > \tau\}$, by the triangle
inequality and $\tau > M \geq \norm{\nabla f(x_t)}$:
\[
  \norm{g_t} - \tau
  \leq \norm{\varepsilon_t} + \norm{\nabla f(x_t)} - \tau
  \leq \norm{\varepsilon_t} + M - \tau
  \leq \norm{\varepsilon_t}.
\]
Moreover, $\norm{g_t} > \tau$ and $\norm{g_t} \leq M + \norm{\varepsilon_t}$
imply $\norm{\varepsilon_t} > \tau - M$. Therefore:
\[
  \bigl(\norm{g_t} - \tau\bigr)^+
  \leq \norm{\varepsilon_t} \cdot \mathbf{1}\!\left(\norm{\varepsilon_t} > \tau - M\right).
\]
Applying Lemma~\ref{lem:pointwise}(a) with $a = \norm{\varepsilon_t}$ and
threshold $\tau - M$:
\[
  \norm{\varepsilon_t} \cdot \mathbf{1}(\norm{\varepsilon_t} > \tau - M)
  \leq \frac{\norm{\varepsilon_t}^p}{(\tau - M)^{p-1}}.
\]
Taking conditional expectations:
\[
  \norm{b_t}
  \leq \E\!\left[\frac{\norm{\varepsilon_t}^p}{(\tau - M)^{p-1}} \;\middle|\; x_t\right]
  = \frac{\E[\norm{\varepsilon_t}^p \mid x_t]}{(\tau - M)^{p-1}}
  \leq \frac{\sigma^p}{(\tau - M)^{p-1}}.
\]
When $\tau \geq 2M$: $(\tau - M) \geq \tau/2$, so $(\tau - M)^{p-1} \geq (\tau/2)^{p-1}
= \tau^{p-1}/2^{p-1}$, giving $\norm{b_t} \leq 2^{p-1}\sigma^p/\tau^{p-1}$.
\end{proof}

\begin{lemma}[Clipped gradient second moment bound]\label{lem:variance}
Under the same setup as Lemma~\ref{lem:bias}:
\[
  \E\!\left[\norm{g_t^c}^2 \mid x_t\right]
  \;\leq\; 2^{p-1}\tau^{2-p}\!\left(M^p + \sigma^p\right).
\]
\end{lemma}

\begin{proof}
By Lemma~\ref{lem:pointwise}(b) with $a = \norm{g_t}$:
\[
  \norm{g_t^c}^2 = \min(\norm{g_t}, \tau)^2 \leq \tau^{2-p}\norm{g_t}^p.
\]
By the triangle inequality $\norm{g_t} \leq \norm{\nabla f(x_t)} + \norm{\varepsilon_t}
\leq M + \norm{\varepsilon_t}$, and the convexity inequality $(a + b)^p \leq
2^{p-1}(a^p + b^p)$ valid for $a, b \geq 0$ and $p \geq 1$:
\[
  \norm{g_t}^p \leq (M + \norm{\varepsilon_t})^p \leq 2^{p-1}(M^p + \norm{\varepsilon_t}^p).
\]
Taking conditional expectations:
\[
  \E[\norm{g_t^c}^2 \mid x_t]
  \leq \tau^{2-p} \cdot 2^{p-1}\!\left(M^p + \E[\norm{\varepsilon_t}^p \mid x_t]\right)
  \leq 2^{p-1}\tau^{2-p}(M^p + \sigma^p).
  \qedhere
\]
\end{proof}

\begin{definition}[Bias and variance constants]\label{def:BV}
For a given threshold $\tau > 2M$, define the bias and variance constants:
\[
  B(\tau) := \frac{2^{p-1}\sigma^p}{\tau^{p-1}}, \qquad
  V(\tau)^2 := 2^{p-1}\tau^{2-p}(M^p + \sigma^p).
\]
By Lemmas~\ref{lem:bias} and~\ref{lem:variance}: $\norm{b_t} \leq B(\tau)$
and $\E[\norm{g_t^c}^2 \mid x_t] \leq V(\tau)^2$ for all $t$.
\end{definition}


% ===========================================================================
\section{Median Scale Estimation}
\label{sec:median}
% ===========================================================================

In the adaptive variant, $\tau_t = C \cdot \shat_t$ with $\shat_t$ the median of
the lagged window of norms (Section~\ref{sec:setup}). The goal here is a deviation bound for $\shat_t$.

\begin{lemma}[Median concentration]\label{lem:median}
Let $X_1, \ldots, X_W$ be independent (but not necessarily identically
distributed) random variables. Suppose:
\begin{enumerate}[label=(\roman*),noitemsep]
  \item \textbf{Common median:} There exists $m \in \R$ such that
    $\Prob(X_i \leq m) \geq 1/2$ and $\Prob(X_i \geq m) \geq 1/2$ for all $i$.
  \item \textbf{Uniform density near median:} There exists $c > 0$ such that
    for all $i$ and all $\delta \in [0, R]$:
    \[
      F_i(m + \delta) - F_i(m) \geq c\delta
      \quad\text{and}\quad
      F_i(m) - F_i(m - \delta) \geq c\delta,
    \]
    where $F_i(x) := \Prob(X_i \leq x)$. Equivalently, each $X_i$ has a
    density (or sub-density) that is bounded below by $c$ in a neighborhood
    of $m$.
\end{enumerate}
Let $\hat{m}_W = \operatorname{median}(X_1, \ldots, X_W)$ be the sample median. Then:
\[
  \Prob\!\left(\abs{\hat{m}_W - m} > \delta\right)
  \leq 2\exp\!\left(-2W(c\delta)^2\right).
\]
In particular, taking $c = 1/R$ for a range parameter $R$ gives
$\Prob(\abs{\hat{m}_W - m} > \delta) \leq 2\exp(-2W\delta^2/R^2)$.
\end{lemma}

\begin{proof}
The sample median $\hat{m}_W$ exceeds $m + \delta$ only if fewer than half
of the observations fall below $m + \delta$. Define $Y_i = \mathbf{1}(X_i \leq m + \delta)$.
Since $\Prob(X_i \leq m + \delta) \geq 1/2 + c\delta$ (by the Lipschitz condition),
we have $\E[Y_i] \geq 1/2 + c\delta$. The $Y_i$ are independent Bernoulli, and
$\hat{m}_W > m + \delta$ implies $\sum_{i=1}^W Y_i < W/2$.

By Hoeffding's inequality for the sum of independent bounded random variables:
\[
  \Prob\!\left(\sum_{i=1}^W Y_i < \frac{W}{2}\right)
  \leq \Prob\!\left(\sum_{i=1}^W Y_i < W\E[Y_1] - Wc\delta\right)
  \leq \exp\!\left(-2W(c\delta)^2\right).
\]
A symmetric argument applies for $\hat{m}_W < m - \delta$. A union bound
gives $\Prob(\abs{\hat{m}_W - m} > \delta) \leq 2\exp(-2W(c\delta)^2)$.
\end{proof}

\begin{lemma}[Median sandwich for non-i.d.\ variables]\label{lem:median_sandwich}
Let $X_1, \ldots, X_W$ be independent random variables (not necessarily
identically distributed). Suppose there exist $m_0 \in \R$, $M_0 \geq 0$,
and $c > 0$ such that for every $i$ and every $\delta \in [0, R]$:
\begin{enumerate}[label=(\roman*),noitemsep]
  \item \emph{Approximate-median band.} The population median $m_i$ of $X_i$
    satisfies $\abs{m_i - m_0} \leq M_0$.
  \item \emph{Local slope.} For every $i$ and every
    $y \in [m_0 - M_0 - R, \, m_0 + M_0 + R]$ and every $\delta \in [0, R]$,
    \[
      F_i(y + \delta) - F_i(y) \geq c\,\delta.
    \]
\end{enumerate}
Let $\hat{m}_W = \operatorname{median}(X_1, \ldots, X_W)$ be the sample median.
Then for $0 \leq \delta \leq R$:
\begin{equation}\label{eq:median_sandwich}
  \Prob\!\left(\hat{m}_W > m_0 + M_0 + \delta\right) \leq \exp(-2W(c\delta)^2),
  \qquad
  \Prob\!\left(\hat{m}_W < m_0 - M_0 - \delta\right) \leq \exp(-2W(c\delta)^2).
\end{equation}
In particular, $\Prob(\abs{\hat{m}_W - m_0} > M_0 + \delta) \leq 2\exp(-2W(c\delta)^2)$.
\end{lemma}

\begin{proof}
Fix $\delta \in [0, R]$ and put $y := m_0 + M_0 + \delta$. By condition~(i),
$m_i \leq m_0 + M_0$, so by condition~(ii) applied at the base point
$m_0 + M_0$:
\[
  \Prob(X_i \leq y)
  \;=\; F_i(m_0 + M_0 + \delta)
  \;\geq\; F_i(m_i) + (F_i(m_0 + M_0) - F_i(m_i)) + c\delta
  \;\geq\; \tfrac{1}{2} + c\delta,
\]
where the last step uses $F_i(m_i) \geq 1/2$ and the monotonicity of $F_i$.
Hence $Y_i := \mathbf 1(X_i \leq y)$ are independent Bernoullis with mean
$\geq 1/2 + c\delta$. The event $\{\hat m_W > y\}$ implies $\sum_i Y_i < W/2$;
Hoeffding's inequality applied to $\{1 - Y_i\}$ (so the deviation is at least
$Wc\delta$) gives
\[
  \Prob(\hat m_W > m_0 + M_0 + \delta)
  \;\leq\; \exp(-2W(c\delta)^2).
\]
The lower-tail bound is symmetric: take $y = m_0 - M_0 - \delta$, use
$m_i \geq m_0 - M_0$ and $F_i(m_i) \leq 1/2$ to get
$\Prob(X_i \geq y) \geq 1/2 + c\delta$, then apply Hoeffding.
\end{proof}

\begin{remark}[Why a shift identity is \emph{not} used]\label{rem:no_shift}
For $d > 1$, the distribution of $\norm{\nabla f(x_s) + \varepsilon_s}$ is
\emph{not} a horizontal shift of the distribution of $\norm{\varepsilon_s}$:
adding a deterministic vector $b\in\R^d$ to the random vector $\varepsilon_s$
changes the law of $\norm{\,\cdot\,}$ in a direction-dependent way, not a
location-family fashion. We avoid the shift identity entirely;
Lemma~\ref{lem:median_sandwich} only uses the deterministic sandwich
$\norm{\varepsilon_s}-M\leq\norm{g_s}\leq\norm{\varepsilon_s}+M$, which is
correct in every dimension by the reverse triangle inequality.
\end{remark}

\begin{lemma}[Verification of the sandwich hypothesis for $\norm{g_s}$]\label{lem:conditioning}
Under Assumptions~\ref{ass:heavy} (conditional independence) and~\ref{ass:signal},
along the lagged-window indices $s \in [\max(1, t-W),\, t-1]$ used in the
algorithm:
\begin{enumerate}[label=\textup{(}\alph*\textup{)},noitemsep,leftmargin=*]
  \item Conditioned on the trajectory $\{x_s\}_{s\leq t-1}$, the gradient norms
    $\{\norm{g_s}\}$ in the window are independent.
  \item Each $\norm{g_s}$ has population median $m_s$ satisfying
    $\abs{m_s - \sigma_{\mathrm{med}}} \leq M$, where
    $\sigma_{\mathrm{med}}$ is the conditional median of $\norm{\varepsilon_s}$
    given $x_s$ (assumed to be the same constant for all $s$ on the trajectory;
    see Remark~\ref{rem:sigma_med_const}).
  \item If, in addition, the conditional distribution of $\norm{\varepsilon_s}$
    given $x_s$ admits a density bounded below by $c_{\mathrm{cdf}}>0$
    on the slab $[\sigma_{\mathrm{med}} - M - R,\, \sigma_{\mathrm{med}} + M + R]$
    uniformly over $x_s$, then condition~(ii) of
    Lemma~\ref{lem:median_sandwich} holds with $c = c_{\mathrm{cdf}}$ and the
    same $R$, after the substitution $X_s \mapsto \norm{g_s}$,
    $m_0 \mapsto \sigma_{\mathrm{med}}$, $M_0 \mapsto M$.
\end{enumerate}
\end{lemma}

\begin{proof}
\textbf{(a)} By Assumption~\ref{ass:heavy}(c), $\{\varepsilon_s\}$ are
conditionally independent given the trajectory; so are
$\norm{g_s} = h(x_s,\varepsilon_s)$ for the deterministic function
$h(x,e):=\norm{\nabla f(x)+e}$.

\textbf{(b)} The reverse triangle inequality gives the deterministic sandwich
$\norm{\varepsilon_s} - M \leq \norm{g_s} \leq \norm{\varepsilon_s} + M$ for
each $s$ (using $\norm{\nabla f(x_s)} \leq M$). Hence the $1/2$-quantile of
$\norm{g_s}$ lies in $[\sigma_{\mathrm{med}} - M, \sigma_{\mathrm{med}} + M]$.

\textbf{(c)} The slab-density hypothesis is imposed \emph{directly} on the
gradient norm law: we assume that conditional on $x_s$, the law of
$\norm{g_s}$ admits a density (or sub-density) $f_s(\cdot\mid x_s)$ bounded
below by $c_{\mathrm{cdf}}$ on the closed slab
$\mathcal S := [\sigma_{\mathrm{med}} - M - R,\, \sigma_{\mathrm{med}} + M + R]$,
uniformly over $x_s$. Sufficient conditions are listed in
Remark~\ref{rem:slab_alternatives}. Given the hypothesis, condition~(ii) of
Lemma~\ref{lem:median_sandwich} holds at every base point in $\mathcal S$:
for $y \in \mathcal S$ and $\delta \in [0, R]$,
$F_{\norm{g_s}}(y + \delta) - F_{\norm{g_s}}(y) \geq c_{\mathrm{cdf}}\,\delta$
by integration, after the substitution
$X_s \mapsto \norm{g_s}$, $m_0 \mapsto \sigma_{\mathrm{med}}$,
$M_0 \mapsto M$, $c \mapsto c_{\mathrm{cdf}}$.
\end{proof}

\begin{remark}[On Assumption~\ref{ass:tail} alone being insufficient]\label{rem:tail_not_density}
A regularly-varying \emph{tail} bound (Assumption~\ref{ass:tail}) does
\emph{not} by itself produce a density lower bound near the median: it only
controls the mass beyond a level set. Lemma~\ref{lem:conditioning}(c)
therefore augments Assumption~\ref{ass:tail} with an explicit slab density
hypothesis on the noise norm. In our Phase~1 experiments,
$\varepsilon_s\sim$ Pareto with deterministic radial density satisfies this
hypothesis with explicit $(c_{\mathrm{cdf}}, R)$. We discuss alternative
sufficient conditions in Remark~\ref{rem:slab_alternatives}.
\end{remark}

\begin{remark}[$\sigma_{\mathrm{med}}$ as a constant of the problem]\label{rem:sigma_med_const}
Throughout, $\sigma_{\mathrm{med}}$ denotes the conditional median of
$\norm{\varepsilon_t}$ given $x_t$, which we assume is the same constant for
all $t$. This holds automatically when the noise $\varepsilon_t$ has a
distribution that does not depend on $x_t$ (Phase~1 setting). When the noise
distribution depends on $x_t$ but Assumption~\ref{ass:tail} provides a uniform
control, one can replace $\sigma_{\mathrm{med}}$ by the supremum
$\overline\sigma_{\mathrm{med}} := \sup_{x}\operatorname{med}(\norm{\varepsilon_t}\mid x_t = x)$
and an analogous infimum, and absorb the gap into the constant $M_0$ of
Lemma~\ref{lem:median_sandwich}.
\end{remark}

\begin{remark}[Alternative sufficient conditions for slab density]\label{rem:slab_alternatives}
Two cleaner sufficient conditions for the slab-density hypothesis used in
Lemma~\ref{lem:conditioning}(c):
\begin{itemize}[noitemsep,leftmargin=*]
  \item \emph{Bounded direction.} If $\nabla f(x_s) = 0$ on the trajectory
    (e.g., centered iterates), then $\norm{g_s} = \norm{\varepsilon_s}$ and
    Assumption~\ref{ass:tail} can be \emph{strengthened} to a density lower
    bound by adding a $C^1$ regularity condition on the noise law.
  \item \emph{Pareto-spheroidal noise.} If $\varepsilon_s = R_s\,U_s$ with
    $R_s\sim$ Pareto and $U_s$ uniform on the sphere, independent, and $R_s$
    has a $C^0$ density bounded below by $c_R$ on a neighborhood of
    $\sigma_{\mathrm{med}}$, then $\norm{\varepsilon_s} = R_s$ has the same
    density and $c_{\mathrm{cdf}} = c_R$ holds with no further work. This
    matches the Phase~1 sandbox.
\end{itemize}
\end{remark}

\begin{proposition}[High-probability scale estimate]\label{prop:scale}
Adopt the lagged-window contract: at step $t$, let
$\shat_t := \operatorname{median}\bigl\{\norm{g_s} : s \in [\max(1, t-W), t-1]\bigr\}$
be the rolling-window median, computed \emph{before} $\norm{g_t}$ enters the
window (so $\shat_t \in \Filt_t$). Under
Assumptions~\ref{ass:heavy}--\ref{ass:signal} and the slab-density hypothesis
of Lemma~\ref{lem:conditioning}\textup{(c)}, with the constant
$c_{\mathrm{cdf}} > 0$ from that hypothesis and any $\delta \in (0, R]$:
\[
  \Prob\!\left(\abs{\shat_t - \sigma_{\mathrm{med}}} > M + \delta\right)
  \;\leq\; 2\exp\!\left(-2W c_{\mathrm{cdf}}^2\delta^2\right).
\]
Equivalently, fixing any $\varepsilon \in (0, 1/2]$ and putting $\delta := \varepsilon \sigma_{\mathrm{med}}$ (assumed $\leq R$), the event
\[
  \mathcal{E}_t := \bigl\{\,\abs{\shat_t - \sigma_{\mathrm{med}}}
    \;\leq\; M + \varepsilon\sigma_{\mathrm{med}}\,\bigr\}
\]
holds with probability at least
$1 - 2\exp\bigl(-2W c_{\mathrm{cdf}}^2\varepsilon^2\sigma_{\mathrm{med}}^2\bigr)$.
Applying a union bound over $t \in [W+1, T]$:
\[
  \Prob\!\left(\bigcap_{t=W+1}^T \mathcal{E}_t\right)
  \;\geq\; 1 - 2T\exp\bigl(-2W c_{\mathrm{cdf}}^2\varepsilon^2\sigma_{\mathrm{med}}^2\bigr).
\]
\end{proposition}

\begin{proof}
By Lemma~\ref{lem:conditioning}\textup{(a)} the window observations are
conditionally independent. By Lemma~\ref{lem:conditioning}\textup{(b)} each
window observation has population median in
$[\sigma_{\mathrm{med}} - M, \sigma_{\mathrm{med}} + M]$. By
Lemma~\ref{lem:conditioning}\textup{(c)} the slab-density hypothesis
verifies condition~(ii) of Lemma~\ref{lem:median_sandwich} on the slab
$[\sigma_{\mathrm{med}} - M - R, \sigma_{\mathrm{med}} + M + R]$ with constant
$c = c_{\mathrm{cdf}}$. Applying Lemma~\ref{lem:median_sandwich} with
$m_0 = \sigma_{\mathrm{med}}$ and $M_0 = M$ yields
\eqref{eq:median_sandwich}, and a union bound over the two tails delivers
the deviation bound. The choice $\delta = \varepsilon\sigma_{\mathrm{med}}$
gives $\mathcal{E}_t$ with the stated probability, and the union bound
over $t \in [W+1, T]$ closes the argument.
\end{proof}

\begin{remark}[$\sigma_{\mathrm{med}} \gg M$ regime]\label{rem:hi_noise}
Theorems~\ref{thm:adaptive} and~\ref{thm:coupled} below use the simplification
$\shat_t \in [(1-\varepsilon')\sigma_{\mathrm{med}}, (1+\varepsilon')\sigma_{\mathrm{med}}]$
with $\varepsilon' = \varepsilon + M/\sigma_{\mathrm{med}}$. This requires
$\sigma_{\mathrm{med}} > 2M$ \emph{(or any explicit lower bound that controls
$\varepsilon'$)}, which we adopt as a problem-instance regime. When this fails,
Proposition~\ref{prop:scale}'s additive $M$ term cannot be absorbed
multiplicatively, and the bias/variance perturbation must be tracked
separately. We make this regime explicit in the statements of the affected
theorems.
\end{remark}


% ===========================================================================
\section{Fixed-Schedule Convergence}
\label{sec:fixed}
% ===========================================================================

\begin{theorem}[Main theorem: fixed schedule]\label{thm:fixed}
Under Assumptions~\ref{ass:smooth}--\ref{ass:signal}, let $\alpha$-Robust SGD
run for $T$ steps with constant learning rate $\eta$ and clipping threshold
$\tau \geq 2M$, where $M$ is from Assumption~\ref{ass:signal}. Then:
\begin{equation}\label{eq:general_bound}
  \frac{1}{T}\sum_{t=1}^{T}\E\!\left[\norm{\nabla f(x_t)}^2\right]
  \;\leq\;
  \frac{2\Delta_f}{\eta T}
  \;+\; B(\tau)^2
  \;+\; L\eta\, V(\tau)^2,
\end{equation}
where $B(\tau)$ and $V(\tau)$ are from Definition~\ref{def:BV}.

\medskip

Setting $\tau$ and $\eta$ optimally as functions of $T$ (the explicit
constants are derived in the proof below; see eq.~\eqref{eq:tau_opt}):
\begin{equation}\label{eq:optimal_rate}
  \tau^*(T) \;=\; \left(\frac{A_B}{2\sqrt{2L\Delta_f\,A_V}}\right)^{\!\frac{2}{3p-2}}
    \cdot T^{\frac{1}{3p-2}},
  \qquad
  \eta^*(T) \;=\; \sqrt{\frac{2\Delta_f}{L\, V(\tau^*)^2\, T}},
\end{equation}
where $A_B := 2^{2(p-1)}\sigma^{2p}$ and $A_V := 2^{p-1}(M^p+\sigma^p)$. The
bound becomes:
\begin{equation}\label{eq:rate}
  \frac{1}{T}\sum_{t=1}^{T}\E\!\left[\norm{\nabla f(x_t)}^2\right]
  \;\leq\;
  C_{\mathrm{rate}} \cdot T^{-\frac{2(p-1)}{3p-2}},
\end{equation}
where $C_{\mathrm{rate}}$ depends on $L, \Delta_f, \sigma, M, p$ but not on $T$.
\end{theorem}

\begin{proof}
\textbf{Step 1: One-step descent.}
By Assumption~\ref{ass:smooth} (the descent lemma for $L$-smooth functions):
\begin{equation}\label{eq:descent_raw}
  f(x_{t+1}) \leq f(x_t) + \inner{\nabla f(x_t)}{x_{t+1} - x_t}
    + \frac{L}{2}\norm{x_{t+1} - x_t}^2.
\end{equation}
Substituting $x_{t+1} - x_t = -\eta\, g_t^c$:
\begin{equation}\label{eq:descent_clip}
  f(x_{t+1}) \leq f(x_t) - \eta\inner{\nabla f(x_t)}{g_t^c}
    + \frac{L\eta^2}{2}\norm{g_t^c}^2.
\end{equation}

\medskip
\textbf{Step 2: Conditional expectation.}
Take $\E[\cdot \mid \Filt_t]$ where $\Filt_t = \sigma(x_1, \varepsilon_1,
\ldots, \varepsilon_{t-1})$ is the filtration making $x_t$ measurable.
Since $g_t^c$ depends on $\varepsilon_t$ (and hence is not $\Filt_t$-measurable),
while $\nabla f(x_t)$ is $\Filt_t$-measurable:
\begin{align}
  \E[f(x_{t+1}) \mid \Filt_t]
  &\leq f(x_t) - \eta\inner{\nabla f(x_t)}{\E[g_t^c \mid \Filt_t]}
    + \frac{L\eta^2}{2}\E\!\left[\norm{g_t^c}^2 \mid \Filt_t\right] \notag\\
  &= f(x_t) - \eta\inner{\nabla f(x_t)}{\nabla f(x_t) - b_t}
    + \frac{L\eta^2}{2}\E\!\left[\norm{g_t^c}^2 \mid \Filt_t\right],
    \label{eq:descent_cond}
\end{align}
where $b_t = \nabla f(x_t) - \E[g_t^c \mid \Filt_t]$ is the clipping bias.

\medskip
\textbf{Step 3: Bounding the inner product via Young's inequality.}
Expanding the inner product:
\[
  \inner{\nabla f(x_t)}{\nabla f(x_t) - b_t}
  = \norm{\nabla f(x_t)}^2 - \inner{\nabla f(x_t)}{b_t}.
\]
By Young's inequality with parameter $1$ (i.e., $\inner{a}{b} \leq
\frac{1}{2}\norm{a}^2 + \frac{1}{2}\norm{b}^2$):
\[
  \inner{\nabla f(x_t)}{b_t}
  \leq \frac{1}{2}\norm{\nabla f(x_t)}^2 + \frac{1}{2}\norm{b_t}^2.
\]
Therefore:
\begin{equation}\label{eq:inner_bound}
  \inner{\nabla f(x_t)}{\nabla f(x_t) - b_t}
  \geq \frac{1}{2}\norm{\nabla f(x_t)}^2 - \frac{1}{2}\norm{b_t}^2.
\end{equation}

\medskip
\textbf{Step 4: Substituting bounds.}
Inserting~\eqref{eq:inner_bound} into~\eqref{eq:descent_cond}:
\begin{equation}\label{eq:descent_bv}
  \E[f(x_{t+1}) \mid \Filt_t]
  \leq f(x_t) - \frac{\eta}{2}\norm{\nabla f(x_t)}^2
    + \frac{\eta}{2}\norm{b_t}^2
    + \frac{L\eta^2}{2}\E\!\left[\norm{g_t^c}^2 \mid \Filt_t\right].
\end{equation}

By Lemma~\ref{lem:bias} and Lemma~\ref{lem:variance}:
\[
  \norm{b_t}^2 \leq B(\tau)^2, \qquad
  \E[\norm{g_t^c}^2 \mid \Filt_t] \leq V(\tau)^2.
\]
Substituting:
\begin{equation}\label{eq:descent_final}
  \E[f(x_{t+1}) \mid \Filt_t]
  \leq f(x_t) - \frac{\eta}{2}\norm{\nabla f(x_t)}^2
    + \frac{\eta}{2}B(\tau)^2
    + \frac{L\eta^2}{2}V(\tau)^2.
\end{equation}

\medskip
\textbf{Step 5: Telescoping.}
Taking full expectations in~\eqref{eq:descent_final} and summing from $t = 1$ to $T$:
\[
  \E[f(x_{T+1})] - f(x_1) \leq
    -\frac{\eta}{2}\sum_{t=1}^T \E\!\left[\norm{\nabla f(x_t)}^2\right]
    + \frac{T\eta}{2}B(\tau)^2 + \frac{TL\eta^2}{2}V(\tau)^2.
\]
Rearranging and using $\E[f(x_{T+1})] \geq f^*$:
\[
  \frac{\eta}{2}\sum_{t=1}^T \E\!\left[\norm{\nabla f(x_t)}^2\right]
  \leq (f(x_1) - f^*) + \frac{T\eta}{2}B(\tau)^2 + \frac{TL\eta^2}{2}V(\tau)^2.
\]
Dividing by $T\eta/2$:
\begin{equation}\label{eq:prelim_bound}
  \frac{1}{T}\sum_{t=1}^T \E\!\left[\norm{\nabla f(x_t)}^2\right]
  \leq \underbrace{\frac{2\Delta_f}{\eta T}}_{\text{(I)}}
    + \underbrace{B(\tau)^2}_{\text{(II)}}
    + \underbrace{L\eta\, V(\tau)^2}_{\text{(III)}}.
\end{equation}
This establishes~\eqref{eq:general_bound}.

\medskip
\textbf{Step 6: Optimizing $\eta$ (balancing terms I and III).}
Terms~(I) and~(III) have opposite monotonicity in $\eta$:
(I)~$= 2\Delta_f/(\eta T)$ decreases in $\eta$, while
(III)~$= L\eta V^2$ increases. Their sum is minimized when they are equal:
\[
  \frac{2\Delta_f}{\eta T} = L\eta V^2
  \quad\Longrightarrow\quad
  \eta^* = \sqrt{\frac{2\Delta_f}{L T V^2}}.
\]
Substituting $\eta^*$ back:
\[
  \text{(I) + (III)} = 2\sqrt{\frac{2L\Delta_f V^2}{T}}.
\]
The bound becomes:
\begin{equation}\label{eq:after_eta}
  \frac{1}{T}\sum \E\!\left[\norm{\nabla f(x_t)}^2\right]
  \leq 2\sqrt{\frac{2L\Delta_f\, V(\tau)^2}{T}} + B(\tau)^2.
\end{equation}

\medskip
\textbf{Step 7: Optimizing $\tau$ (balancing the two remaining terms).}
Write $B(\tau)^2$ and $V(\tau)^2$ in terms of $\tau$:
\[
  B(\tau)^2 = \frac{2^{2(p-1)}\sigma^{2p}}{\tau^{2(p-1)}}
  =: \frac{A_B}{\tau^{2(p-1)}},
\]
\[
  V(\tau)^2 = 2^{p-1}(M^p + \sigma^p)\,\tau^{2-p}
  =: A_V \cdot \tau^{2-p},
\]
where $A_B := 2^{2(p-1)}\sigma^{2p}$ and $A_V := 2^{p-1}(M^p + \sigma^p)$.

Substituting into~\eqref{eq:after_eta}:
\begin{equation}\label{eq:two_terms}
  \text{bound} \leq
  \underbrace{2\sqrt{\frac{2L\Delta_f A_V}{T}} \cdot \tau^{(2-p)/2}}_{\text{Term A}}
  + \underbrace{\frac{A_B}{\tau^{2(p-1)}}}_{\text{Term B}}.
\end{equation}
Term~A increases in $\tau$ (since $(2-p)/2 > 0$) while Term~B decreases.
Setting them equal:
\[
  2\sqrt{\frac{2L\Delta_f A_V}{T}} \cdot \tau^{(2-p)/2}
  = \frac{A_B}{\tau^{2(p-1)}}.
\]
Solving for $\tau$:
\[
  \tau^{(2-p)/2 + 2(p-1)} = \frac{A_B}{2\sqrt{2L\Delta_f A_V/T}}
  = \frac{A_B \sqrt{T}}{2\sqrt{2L\Delta_f A_V}}.
\]
The exponent on $\tau$ is:
\[
  \frac{2-p}{2} + 2(p-1) = \frac{2 - p + 4p - 4}{2} = \frac{3p - 2}{2}.
\]
Therefore:
\begin{equation}\label{eq:tau_opt}
  \tau^{(3p-2)/2} = \frac{A_B\sqrt{T}}{2\sqrt{2L\Delta_f A_V}}
  \quad\Longrightarrow\quad
  \tau^* = \left(\frac{A_B}{2\sqrt{2L\Delta_f A_V}}\right)^{\frac{2}{3p-2}}
    T^{\frac{1}{3p-2}}.
\end{equation}
The key dependence on $T$ is $\tau^* \propto T^{1/(3p-2)}$.

\medskip
\textbf{Step 8: Final rate.}
Substituting $\tau^*$ back into Term~B of~\eqref{eq:two_terms}:
\[
  B(\tau^*)^2 = \frac{A_B}{(\tau^*)^{2(p-1)}} \propto T^{-2(p-1)/(3p-2)}.
\]
Explicitly:
\[
  (\tau^*)^{2(p-1)} = \left(\frac{A_B}{2\sqrt{2L\Delta_f A_V}}\right)^{\frac{4(p-1)}{3p-2}}
    T^{\frac{2(p-1)}{3p-2}},
\]
so:
\begin{equation}
  B(\tau^*)^2 = A_B^{1 - \frac{4(p-1)}{3p-2}} \cdot
    \left(2\sqrt{2L\Delta_f A_V}\right)^{\frac{4(p-1)}{3p-2}}
    \cdot T^{-\frac{2(p-1)}{3p-2}}.
\end{equation}
Since Terms~A and~B are equal at $\tau = \tau^*$, the total bound is:
\begin{equation}\label{eq:final_rate}
  \frac{1}{T}\sum_{t=1}^{T}\E\!\left[\norm{\nabla f(x_t)}^2\right]
  \leq C_{\mathrm{rate}} \cdot T^{-\frac{2(p-1)}{3p-2}},
\end{equation}
where $C_{\mathrm{rate}}$ is an explicit constant depending only on
$L, \Delta_f, \sigma, M, p$:
\[
  C_{\mathrm{rate}} = 2 A_B^{\frac{3p - 2 - 4(p-1)}{3p-2}}
    \left(2\sqrt{2L\Delta_f A_V}\right)^{\frac{4(p-1)}{3p-2}}.
\]

\medskip
\textbf{Step 9: Verify the optimal $\eta$.}
With $\tau^* \propto T^{1/(3p-2)}$, we have $V(\tau^*)^2 = A_V (\tau^*)^{2-p}
\propto T^{(2-p)/(3p-2)}$. The optimal learning rate is:
\[
  \eta^* = \sqrt{\frac{2\Delta_f}{LT V(\tau^*)^2}}
  \propto T^{-1/2 - (2-p)/(2(3p-2))}
  = T^{-\frac{3p-2+2-p}{2(3p-2)}}
  = T^{-\frac{2p}{2(3p-2)}}
  = T^{-\frac{p}{3p-2}}.
\]
For $p = 2$: $\eta^* \propto T^{-1/2}$ (classical). For $p = 1.5$:
$\eta^* \propto T^{-3/5}$.

\medskip
\textbf{Step 10: Verify $\tau^* \geq 2M$ for large $T$.}
We require $\tau^* \geq 2M$ for Lemma~\ref{lem:bias} to apply.
Since $\tau^* \propto T^{1/(3p-2)} \to \infty$ and $M$ is a constant,
this holds for all $T \geq T_0$ with $T_0 = (2M)^{3p-2} \cdot (\text{const})$.
For $T < T_0$, one can use a fixed threshold $\tau_0 \geq 2M$ (the burn-in
period), incurring $O(T_0) = O(1)$ additive cost.
\end{proof}

\begin{remark}[Rate interpolation]\label{rem:interp}
The rate $T^{-2(p-1)/(3p-2)}$ interpolates correctly:
\begin{center}
\renewcommand{\arraystretch}{1.2}
\begin{tabular}{lll}
\toprule
$p$ & Rate exponent & Interpretation \\
\midrule
$2$ & $-2/4 = -1/2$ & Classical non-convex SGD with bounded variance \\
$1.5$ & $-1/2.5 = -2/5$ & Intermediate heavy tails \\
$1.2$ & $-0.4/1.6 = -1/4$ & Very heavy tails \\
$\to 1^+$ & $\to 0$ & No convergence (first moment barely exists) \\
\bottomrule
\end{tabular}
\end{center}
\end{remark}

\begin{remark}[Known-horizon requirement]\label{rem:horizon}
The optimal $\tau^*$ and $\eta^*$ depend on $T$, requiring the time horizon
to be known in advance. This is standard in the non-convex optimization
literature (cf.\ Zhang et al., 2020). For the unknown-horizon case, one can
apply the \textbf{doubling trick}: run the algorithm in epochs of length
$T_k = 2^k$, restarting with fresh $\tau^*_{T_k}$ and $\eta^*_{T_k}$ at each
epoch. This incurs at most an $O(\log T)$ multiplicative overhead.
\end{remark}

\begin{remark}[Fixed $\tau$ independent of $T$]\label{rem:fixed_tau}
If $\tau$ is set to a fixed constant (not depending on $T$), the bias term
$B(\tau)^2$ is a constant \emph{floor} that does not vanish with $T$.
The bound becomes:
\[
  \frac{1}{T}\sum \E[\norm{\nabla f(x_t)}^2] \leq
  2\sqrt{\frac{2L\Delta_f V(\tau)^2}{T}} + B(\tau)^2.
\]
This converges to $B(\tau)^2$ as $T \to \infty$ — the algorithm reaches a
neighborhood of stationarity but does not converge to zero. In practice,
with $\tau = 5\hat{\sigma}$, the bias floor $B(\tau)^2 \leq 2\sigma^2/5^{2(p-1)}$
is small ($\leq 0.16\sigma^2$ for $p = 1.5$), so the algorithm performs well
for practical time horizons.
\end{remark}


% ===========================================================================
\section{Adaptive Threshold with Scale Estimation}
\label{sec:adaptive}
% ===========================================================================

Consider $\tau_t = C \cdot \shat_t$ with $\shat_t$ from the lagged median window.

\begin{theorem}[Adaptive fixed schedule]\label{thm:adaptive}
Under Assumptions~\ref{ass:smooth}--\ref{ass:signal} and the slab-density
hypothesis of Lemma~\ref{lem:conditioning}\textup{(c)}, fix $C > 0$, $W \geq 1$,
and $\varepsilon \in (0, 1/2]$. Assume the regime $\sigma_{\mathrm{med}} > 2M$
\textup{(}so that $\varepsilon' := \varepsilon + M/\sigma_{\mathrm{med}} \in (0, 1)$\textup{)}
and $C(1-\varepsilon')\sigma_{\mathrm{med}} \geq 2M$. Run $\alpha$-Robust SGD
with $\tau_t = C \cdot \shat_t$ for $t > W$ \textup{(}lagged window;
Remark~\ref{rem:lag}\textup{)} and a fixed $\tau_0 \geq 2M$ for $t \leq W$,
with constant learning rate $\eta$. Define
$\mathcal G := \bigcap_{t=W+1}^T \mathcal E_t$ from Proposition~\ref{prop:scale};
then $\Prob(\mathcal G) \geq 1 - 2T\exp(-2W c_{\mathrm{cdf}}^2 \varepsilon^2 \sigma_{\mathrm{med}}^2)$.

\medskip
\textbf{In-expectation guarantee on $\mathcal{G}$.}
\begin{equation}\label{eq:adaptive_bound}
  \frac{1}{T}\sum_{t=1}^T \E\!\left[\norm{\nabla f(x_t)}^2 \;\Big|\; \mathcal G\right]
  \;\leq\;
  \frac{2\Delta_f}{\eta T}
  + (1 + \delta_W)^2 B(C\sigma_{\mathrm{med}})^2
  + (1 + \delta_W) L\eta\, V(C\sigma_{\mathrm{med}})^2
  + O\!\left(\frac{W}{T}\right),
\end{equation}
where $\delta_W := \max\bigl(4(p-1)\varepsilon',\, 2(2-p)\varepsilon'\bigr) \leq 4\varepsilon'$
captures the estimation overhead.

\medskip
\textbf{Asymptotic order of $\delta_W$.} Choosing
$\varepsilon = \Theta\!\bigl((W c_{\mathrm{cdf}}^2 \sigma_{\mathrm{med}}^2)^{-1/2}\sqrt{\log T}\bigr)$
makes the failure probability $\Prob(\mathcal G^c) = \mathcal{O}(T^{-1})$ and
$\delta_W = \mathcal O(\sqrt{\log T/W} + M/\sigma_{\mathrm{med}})$. The
$M/\sigma_{\mathrm{med}}$ term is a problem constant from the median-vs-norm
sandwich (Remark~\ref{rem:hi_noise}); it does not vanish with $T$ and shows up
only as a multiplicative inflation of the bias/variance constants in the rate
analysis, not in the rate exponent.
\end{theorem}

\begin{proof}
\textbf{Step~1: Lagged-window measurability.}
By the lagged-window contract (Remark~\ref{rem:lag}), the window
$\{\norm{g_s}\}_{s = t-W}^{t-1}$ used to compute $\shat_t$ is
$\Filt_t$-measurable; hence $\tau_t = C\shat_t \in \Filt_t$. Consequently
$\E[g_t^c \mid \Filt_t]$ is computed at fixed $\tau_t$, and
Lemmas~\ref{lem:bias}--\ref{lem:variance} yield, on $\Filt_t$,
\[
  \norm{b_t} \leq B(\tau_t),
  \qquad
  \E[\norm{g_t^c}^2 \mid \Filt_t] \leq V(\tau_t)^2.
\]

\textbf{Step~2: Burn-in ($t \leq W$).}
With $\tau_0 \geq 2M$ constant, eq.~\eqref{eq:descent_final} applies and
yields total cost $f(x_1) - \E[f(x_{W+1})] + \frac{W\eta}{2}B(\tau_0)^2
+ \frac{WL\eta^2}{2}V(\tau_0)^2 = O(W)$, contributing $O(W/T)$ to the
final bound after dividing by $T\eta/2$.

\textbf{Step~3: Threshold envelope on $\mathcal E_t$ ($t > W$).}
On $\mathcal E_t$, $\abs{\shat_t - \sigma_{\mathrm{med}}} \leq M + \varepsilon\sigma_{\mathrm{med}}$,
i.e.\ $\shat_t \in [(1-\varepsilon')\sigma_{\mathrm{med}}, (1+\varepsilon')\sigma_{\mathrm{med}}]$
with $\varepsilon' = \varepsilon + M/\sigma_{\mathrm{med}} \in (0,1)$ by the
regime hypothesis. Therefore
$\tau_t \in [C(1-\varepsilon')\sigma_{\mathrm{med}}, C(1+\varepsilon')\sigma_{\mathrm{med}}]$,
and the lower envelope $\tau_t \geq C(1-\varepsilon')\sigma_{\mathrm{med}} \geq 2M$
\textup{(}by hypothesis\textup{)} ensures the simplified bias bound applies.

\textbf{Step~4: Inflated bias and variance bounds on $\mathcal E_t$.}
Using $B$ decreasing and $V^2$ increasing in $\tau$:
\begin{align*}
  B(\tau_t) &\leq (1-\varepsilon')^{-(p-1)}\,B(C\sigma_{\mathrm{med}})
    \;\leq\; \bigl(1 + 4(p-1)\varepsilon'\bigr)\,B(C\sigma_{\mathrm{med}}), \\
  V(\tau_t)^2 &\leq (1+\varepsilon')^{2-p}\,V(C\sigma_{\mathrm{med}})^2
    \;\leq\; \bigl(1 + 2(2-p)\varepsilon'\bigr)\,V(C\sigma_{\mathrm{med}})^2,
\end{align*}
using the elementary bounds derived in Theorem~\ref{thm:coupled} Steps~3--4
\textup{(}same algebra: $(1-x)^{-(p-1)} \leq 1+4(p-1)x$ and $(1+x)^{2-p}\leq 1+2(2-p)x$
for $x \in [0,1/2]$\textup{)}. Setting
$\delta_W = \max(4(p-1)\varepsilon', 2(2-p)\varepsilon')$ unifies both
factors into $1 + \delta_W$.

\textbf{Step~5: Telescope on $\mathcal G$.}
For $t > W$ on $\mathcal G \subseteq \mathcal E_t$, the descent inequality
$\eqref{eq:descent_final}$ with $\tau$ replaced by $\tau_t$ gives, after
applying the inflated bounds of Step~4 and conditioning by the tower law on
$\Filt_t \cap \mathcal G$:
\[
  \E\!\left[f(x_{t+1}) \mid \Filt_t, \mathcal G\right]
  \leq f(x_t) - \frac{\eta}{2}\norm{\nabla f(x_t)}^2
    + \frac{\eta}{2}(1+\delta_W)^2 B(C\sigma_{\mathrm{med}})^2
    + \frac{L\eta^2}{2}(1+\delta_W) V(C\sigma_{\mathrm{med}})^2.
\]
Telescoping from $t = W+1$ to $T$ on $\mathcal G$, using
$\E[f(x_{T+1})\mid\mathcal G] \geq f^*$, dividing by $T\eta/2$, and adding the
burn-in cost (Step~2) yields~\eqref{eq:adaptive_bound}.

\textbf{Step~6: Filtration--global-event coupling.}
The conditioning $\mathbb E[\,\cdot\,\mid\,\Filt_t,\mathcal G]$ at each step
is well defined because $\mathcal G = \bigcap_{s \geq W+1}^T \mathcal E_s$ is
$\Filt_T$-measurable, so a single regular conditional probability under
$\Prob(\,\cdot\,\mid\mathcal G)$ exists. The tower property
$\E[\E[X \mid \Filt_t,\mathcal G] \mid \mathcal G] = \E[X \mid \mathcal G]$
holds for any integrable random variable $X$, which is what the telescope
uses at each step.
\end{proof}


% ===========================================================================
\section{Hill Estimator Concentration}
\label{sec:hill}
% ===========================================================================

The Hill estimator is used to estimate the tail index $\alpha$ (and hence
$p = \alpha - \varepsilon$ for a safety margin $\varepsilon > 0$). Under
the exact Pareto model, we establish a rigorous finite-sample concentration bound.

\begin{definition}[Hill estimator]\label{def:hill}
Given observations $X_1, \ldots, X_n$ with order statistics
$X_{(1)} \leq \cdots \leq X_{(n)}$, the Hill estimator using the top $k$
order statistics is:
\[
  \ahat_{\mathrm{Hill}}^{-1}
  := \frac{1}{k}\sum_{i=1}^{k}\log\frac{X_{(n-i+1)}}{X_{(n-k)}}.
\]
\end{definition}

\begin{lemma}[R\'enyi representation for exact Pareto]\label{lem:renyi}
Let $X_1, \ldots, X_n$ be i.i.d.\ $\mathrm{Pareto}(\alpha, x_m)$, i.e.,
$\Prob(X > x) = (x_m/x)^\alpha$ for $x \geq x_m > 0$. Then the Hill
estimator satisfies:
\[
  \ahat_{\mathrm{Hill}}^{-1} \stackrel{d}{=} \frac{1}{k}\sum_{j=1}^{k} E_j,
\]
where $E_1, \ldots, E_k$ are i.i.d.\ $\mathrm{Exp}(\alpha)$ (rate $\alpha$,
mean $1/\alpha$).
\end{lemma}

\begin{proof}
By the R\'enyi representation theorem for order statistics: if $U_1, \ldots, U_n$
are i.i.d.\ $\mathrm{Uniform}(0,1)$ with order statistics $U_{(1)} \leq \cdots \leq U_{(n)}$,
then the \emph{spacings} $D_i := U_{(i)} - U_{(i-1)}$ (with $U_{(0)} := 0$)
satisfy the joint distributional identity
$\{D_1, \ldots, D_n\} \stackrel{d}{=} \{G_1/S, \ldots, G_n/S\}$
where $G_i$ are i.i.d.\ $\mathrm{Exp}(1)$ and $S = G_1 + \cdots + G_{n+1}$.

For Pareto random variables, $\log(X_i/x_m)$ are i.i.d.\ $\mathrm{Exp}(\alpha)$.
The log-spacings of the order statistics of exponential random variables are
independent, with $\log(X_{(n-i+1)}/x_m) - \log(X_{(n-i)}/x_m) \sim \mathrm{Exp}(\alpha \cdot i)$.

The Hill estimator's numerator, $\sum_{i=1}^k \log(X_{(n-i+1)}/X_{(n-k)})$,
can be expressed as a linear combination of these independent exponential spacings.
By a classical result (see, e.g., Beirlant et al., 2004, Theorem 4.2.1),
this sum is distributionally equal to a sum of $k$ i.i.d.\ $\mathrm{Exp}(\alpha)$
random variables.
\end{proof}

\begin{lemma}[Finite-sample Hill concentration via Bernstein]\label{lem:hill_conc}
Under the exact Pareto model (Lemma~\ref{lem:renyi}), for any $\varepsilon > 0$
with $\varepsilon \leq \alpha$:
\begin{equation}\label{eq:hill_bound}
  \Prob\!\left(\abs{\ahat_{\mathrm{Hill}} - \alpha} > \varepsilon\right)
  \leq 2\exp\!\left(-\frac{k\varepsilon^2}{16\alpha^2}\right).
\end{equation}
\end{lemma}

\begin{proof}
\textbf{Step 1: Concentration of $\ahat^{-1}$.}
By Lemma~\ref{lem:renyi}, $\ahat^{-1} = \bar{E}$ where
$\bar{E} = \frac{1}{k}\sum_{j=1}^k E_j$ with $E_j \sim \mathrm{Exp}(\alpha)$ i.i.d.

Each $E_j$ has mean $1/\alpha$, variance $1/\alpha^2$, and is sub-exponential
with parameters $(\nu, b) = (1/\alpha, 1/\alpha)$. Specifically, for
$\lambda \in [0, \alpha)$:
\[
  \E[e^{\lambda(E_j - 1/\alpha)}]
  = \frac{\alpha}{\alpha - \lambda}\, e^{-\lambda/\alpha}
  \leq \exp\!\left(\frac{\lambda^2}{2(\alpha - \lambda)\alpha}\right)
  \leq \exp\!\left(\frac{\lambda^2}{\alpha^2}\right)
  \quad\text{for } \lambda \leq \alpha/2.
\]

By Bernstein's inequality for sub-exponential random variables
(see, e.g., Vershynin, 2018, Theorem 2.8.1):
\begin{equation}\label{eq:bernstein}
  \Prob\!\left(\abs{\bar{E} - 1/\alpha} > t\right)
  \leq 2\exp\!\left(-k \cdot \min\!\left(\frac{t^2\alpha^2}{4},\;
    \frac{t\alpha}{2}\right)\right).
\end{equation}

For $t \leq 1/(2\alpha)$ (small deviation regime), the minimum is achieved
by the first term:
\begin{equation}\label{eq:bernstein_small}
  \Prob\!\left(\abs{\bar{E} - 1/\alpha} > t\right)
  \leq 2\exp\!\left(-\frac{kt^2\alpha^2}{4}\right).
\end{equation}

\textbf{Step 2: Translating to $\ahat$ via the delta method.}
The map $\varphi(x) = 1/x$ transforms $\bar{E}$ to $\ahat = 1/\bar{E}$.
On the event $\{\abs{\bar{E} - 1/\alpha} \leq t\}$ with $t < 1/(2\alpha)$:
\[
  \bar{E} \in \left[\frac{1}{\alpha} - t,\; \frac{1}{\alpha} + t\right]
  \subset \left[\frac{1}{2\alpha},\; \frac{3}{2\alpha}\right].
\]
Therefore $\bar{E} \geq 1/(2\alpha) > 0$, and:
\[
  \abs{\ahat - \alpha} = \abs{\frac{1}{\bar{E}} - \frac{1}{1/\alpha}}
  = \frac{\abs{\bar{E} - 1/\alpha}}{\bar{E} \cdot (1/\alpha)}
  \leq \frac{t}{(1/(2\alpha)) \cdot (1/\alpha)}
  = 2\alpha^2 t.
\]

\textbf{Step 3: Combining.}
Given target deviation $\varepsilon > 0$, set $t = \varepsilon/(2\alpha^2)$.
The requirement $t \leq 1/(2\alpha)$ becomes $\varepsilon \leq \alpha$.
Then:
\[
  \Prob(\abs{\ahat - \alpha} > \varepsilon)
  \leq \Prob\!\left(\abs{\bar{E} - 1/\alpha} > \frac{\varepsilon}{2\alpha^2}\right)
  \leq 2\exp\!\left(-\frac{k\varepsilon^2}{4 \cdot 4\alpha^2}\right)
  = 2\exp\!\left(-\frac{k\varepsilon^2}{16\alpha^2}\right).
\]
This completes the proof.
\end{proof}

\begin{remark}[Comparison with asymptotic normality]\label{rem:asymptotic}
The asymptotic normality result of Mason (1982) and Haeusler \& Teugels (1985) gives:
\[
  \sqrt{k}\,(\ahat - \alpha) \xrightarrow{d} \mathcal{N}(0,\, \alpha^2)
  \quad\text{as } k \to \infty,\; k/n \to 0.
\]
For $k = 25$, the normal approximation suggests $\mathrm{SD}(\ahat) \approx \alpha/5$,
which is a tighter prediction than the Bernstein bound~\eqref{eq:hill_bound}
but lacks finite-sample validity. Our Bernstein-based bound~\eqref{eq:hill_bound}
is rigorous for all $k$.
\end{remark}

\begin{remark}[Scope of Lemma~\ref{lem:hill_conc}]\label{rem:beyond_pareto}
Lemma~\ref{lem:hill_conc} assumes exact i.i.d.\ Pareto inputs. Under only
regular variation \textup{(}Assumption~\ref{ass:tail}\textup{)}, log-spacings
are asymptotically exponential and a second-order bias $A(n/k)$ appears;
consistency still holds for $k\to\infty$, $k/n\to 0$, but the finite-sample
constant in Lemma~\ref{lem:hill_conc} does not carry over verbatim.

Gradient norms $\norm{g_s}$ are not i.i.d.\ across $s$ because $x_s$ drifts.
Proposition~\ref{prop:hill_noniid} compares Hill on $\norm{g_s}$ to Hill on
clean norms $\norm{\varepsilon_s}$ by a deterministic perturbation bound.
Putting that together with Lemma~\ref{lem:hill_conc} on $\norm{\varepsilon_s}$
is clean only when the injected noise is i.i.d.\ Pareto \textup{(}Phase~1\textup{)}.

For Phase~2/3 we do not assume a named tail law on gradients; the clamping
$\phat_t\in[p_{\min},p_{\max}]$ gives $\abs{\phat_t-p}\leq 0.99$
\textup{(}eq.~\eqref{eq:clamp_bound}\textup{)} regardless. In that regime
Theorem~\ref{thm:coupled} should be read as convergence to a neighborhood of
stationarity, with radius tied to $\mu_{\max}$, not as the sharp adaptive rate.
\end{remark}


% ===========================================================================
\section{Non-i.i.d.\ Gradient Norms: Quantitative Perturbation Analysis}
\label{sec:noniid}
% ===========================================================================

Gradient norms $\norm{g_t}$ are not i.i.d.\ across $t$. The next proposition
bounds how far Hill on $\{\norm{g_s}\}$ deviates from Hill on clean norms.

\begin{lemma}[Tail index inheritance]\label{lem:tail_inherit}
Under Assumptions~\ref{ass:signal} and~\ref{ass:tail}, conditioned on the
trajectory $\{x_t\}$, the gradient norms
$\norm{g_t} = \norm{\nabla f(x_t) + \varepsilon_t}$
are independent with marginals that share the same tail index $\alpha$.

More precisely, for all $t$ and all $x > 2M$:
\[
  \Prob(\norm{\varepsilon_t} > x + M)
  \leq \Prob(\norm{g_t} > x)
  \leq \Prob(\norm{\varepsilon_t} > x - M).
\]
Since $\Prob(\norm{\varepsilon_t} > x) = x^{-\alpha}L(x)$ with slowly varying $L$,
and $(x \pm M)^{-\alpha}/x^{-\alpha} \to 1$ as $x \to \infty$, the gradient
norms $\norm{g_t}$ have the same tail index $\alpha$ as $\norm{\varepsilon_t}$.
\end{lemma}

\begin{proof}
\textbf{Independence.} By Assumption~\ref{ass:heavy}, the noise terms
$\varepsilon_t$ are conditionally independent given the trajectory. Since
each $\norm{g_t} = h(x_t, \varepsilon_t)$ for the measurable function
$h(x, \varepsilon) = \norm{\nabla f(x) + \varepsilon}$, and $x_t$ is
$\Filt_{t-1}$-measurable, the $\norm{g_t}$ are conditionally independent
given $\{x_s\}_{s \geq 1}$.

\textbf{Tail sandwich.} By the reverse triangle inequality:
$\abs{\norm{g_t} - \norm{\varepsilon_t}} \leq \norm{\nabla f(x_t)} \leq M$.
Therefore $\norm{\varepsilon_t} - M \leq \norm{g_t} \leq \norm{\varepsilon_t} + M$,
which for any $x > 0$ gives:
\[
  \{\norm{\varepsilon_t} > x + M\} \subseteq \{\norm{g_t} > x\}
  \subseteq \{\norm{\varepsilon_t} > x - M\}.
\]
The tail bound follows by taking probabilities.

\textbf{Tail index preservation.} For $x > 2M$:
\[
  \frac{\Prob(\norm{g_t} > x)}{x^{-\alpha}L(x)}
  \in \left[
    \frac{(x+M)^{-\alpha}L(x+M)}{x^{-\alpha}L(x)},\;
    \frac{(x-M)^{-\alpha}L(x-M)}{x^{-\alpha}L(x)}
  \right].
\]
Since $L$ is slowly varying, $L(x \pm M)/L(x) \to 1$ as $x \to \infty$
(by the uniform convergence theorem; see Bingham, Goldie \& Teugels, 1987,
Theorem~1.2.1). Also $(x \pm M)^{-\alpha}/x^{-\alpha}
= (1 \pm M/x)^{-\alpha} \to 1$. By the squeeze theorem:
$\Prob(\norm{g_t} > x) \sim x^{-\alpha}L(x)$, so $\norm{g_t}$ has tail
index $\alpha$.
\end{proof}

\begin{proposition}[Quantitative Hill perturbation bound]\label{prop:hill_noniid}
Under Assumptions~\ref{ass:signal} and~\ref{ass:tail}, let
$Y_s := \norm{\varepsilon_s}$ (clean noise norms) and
$X_s := \norm{g_s} = \norm{\nabla f(x_s) + \varepsilon_s}$ (observed
gradient norms), so that $\abs{X_s - Y_s} \leq M$ for all $s$.
Let $H_k^X$ and $H_k^Y$ denote the Hill estimators (Definition~\ref{def:hill})
applied to $\{X_s\}$ and $\{Y_s\}$ respectively, using the top $k$ order
statistics from a window of size $W$. Then:
\begin{equation}\label{eq:hill_perturbation}
  \abs{(H_k^X)^{-1} - (H_k^Y)^{-1}}
  \leq \frac{4M}{Y_{(W-k)} - M}
  + \frac{2m}{k} \cdot \log\frac{Y_{(W)}}{Y_{(W-k)}},
\end{equation}
where $Y_{(j)}$ denotes the $j$-th order statistic of $\{Y_s\}$, and
$m \leq 2W \cdot \alpha M / Y_{(W-k)}$ is the number of indices whose
top-$k$ membership differs between the $X$ and $Y$ orderings.

Consequently:
\begin{equation}\label{eq:hill_total_error}
  \abs{\ahat_{\mathrm{gradient}} - \alpha}
  \leq \underbrace{\abs{\ahat_{\mathrm{clean}} - \alpha}}_{\text{i.i.d.\ estimation error}}
  + \underbrace{O\!\left(\frac{\alpha^2 M}{Y_{(W-k)} - M}\right)}_{\text{perturbation error}},
\end{equation}
where $\ahat_{\mathrm{gradient}} = 1/(H_k^X)^{-1}$ and
$\ahat_{\mathrm{clean}} = 1/(H_k^Y)^{-1}$.
The perturbation error vanishes as $W \to \infty$ since
$Y_{(W-k)} \to \infty$.

For \textbf{finite} $W$, the clamping $\phat_t \in [p_{\min}, p_{\max}]$
provides the deterministic bound:
\begin{equation}\label{eq:clamp_bound}
  \abs{\phat_t - p} \leq \max(p_{\max} - p,\; p - p_{\min}) \leq 0.99
  \quad \text{for all } p \in (1.01, 1.99).
\end{equation}
\end{proposition}

\begin{proof}
We decompose the perturbation into two effects.

\textbf{Effect 1: Log-ratio perturbation (fixed ordering).}
Suppose the top-$k$ indices are the same for both $\{X_s\}$ and $\{Y_s\}$.
Let $s_1, \ldots, s_k$ be the top-$k$ indices (ordered by $Y$-values).
The $j$-th log-ratio in the Hill estimator changes by:
\begin{align*}
  \Delta_j &:= \log\frac{X_{s_j}}{X_{s_k}} - \log\frac{Y_{s_j}}{Y_{s_k}}
  = \log\frac{X_{s_j}}{Y_{s_j}} - \log\frac{X_{s_k}}{Y_{s_k}}.
\end{align*}
Since $\abs{X_s - Y_s} \leq M$ and $Y_s > Y_{(W-k)} > M$ for the top-$k$
values:
\[
  \abs{\log\frac{X_s}{Y_s}} = \abs{\log\left(1 + \frac{X_s - Y_s}{Y_s}\right)}
  \leq \frac{2\abs{X_s - Y_s}}{Y_s}
  \leq \frac{2M}{Y_{(W-k)}},
\]
using $\abs{\log(1+u)} \leq 2\abs{u}$ for $\abs{u} \leq 1/2$
(valid since $M/Y_{(W-k)} < 1/2$ for $Y_{(W-k)} > 2M$).
Each $\abs{\Delta_j} \leq 4M/Y_{(W-k)}$, so:
\[
  \abs{(H_k^X)^{-1} - (H_k^Y)^{-1}}_{\text{same order}}
  \leq \frac{1}{k}\sum_{j=1}^k \abs{\Delta_j}
  \leq \frac{4M}{Y_{(W-k)}}.
\]

\textbf{Effect 2: Reordering (top-$k$ membership changes).}
An index $s$ can switch top-$k$ membership between the $X$ and $Y$ orderings
only if its gradient norm $X_s$ crosses the threshold. Since
$\abs{X_s - Y_s} \leq M$, this requires $Y_s \in [Y_{(W-k)} - M, Y_{(W-k)} + M]$.

The number of such ``boundary'' indices is:
\[
  m := \#\{s : Y_s \in [Y_{(W-k)} - M, Y_{(W-k)} + M]\}.
\]
By the tail probability of the i.i.d.\ distribution $\Prob(Y \in [a, a+2M])
\leq 2M\alpha a^{-(\alpha+1)}$ for exact Pareto, the expected number of
boundary observations is:
\[
  \E[m] \leq W \cdot 2M\alpha \cdot Y_{(W-k)}^{-(\alpha+1)}.
\]
Since $Y_{(W-k)} \approx (W/k)^{1/\alpha}$ (the quantile), this gives
$\E[m] = O(Mk\alpha / Y_{(W-k)})$.

Each reordered index contributes at most
$\log(Y_{(W)}/Y_{(W-k)})/k$ to the Hill estimator (the maximum
log-ratio). The total reordering error is:
\[
  \abs{(H_k^X)^{-1} - (H_k^Y)^{-1}}_{\text{reorder}}
  \leq \frac{2m}{k} \cdot \log\frac{Y_{(W)}}{Y_{(W-k)}}.
\]
For exact Pareto with $k = W/4$: $\log(Y_{(W)}/Y_{(W-k)}) = O(\log W)$,
and $m/k = O(M\alpha/Y_{(W-k)})$, so this is $O(M\alpha\log W / Y_{(W-k)})$.

\textbf{Total perturbation on $\ahat$.}
Combining Effects 1 and 2, and applying the delta method inversion
(as in Lemma~\ref{lem:hill_conc}, Step 2) to translate from $\ahat^{-1}$
to $\ahat$:
if $\abs{(H_k^X)^{-1} - (H_k^Y)^{-1}} \leq \delta$ with
$\delta < 1/(2\alpha)$, then $\abs{\ahat_X - \ahat_Y} \leq 2\alpha^2\delta$.
The bound~\eqref{eq:hill_total_error} follows with the triangle inequality
$\abs{\ahat_X - \alpha} \leq \abs{\ahat_X - \ahat_Y} + \abs{\ahat_Y - \alpha}$.

The clamping bound~\eqref{eq:clamp_bound} is immediate from
$\phat_t = \operatorname{clamp}(\ahat - \varepsilon, p_{\min}, p_{\max})$.
\end{proof}

\begin{remark}[Convergence of perturbation error]\label{rem:perturbation_rate}
For exact Pareto noise with $\alpha \in (1, 2]$ and $k = W/4$:
$Y_{(W-k)} \approx 4^{1/\alpha} \cdot x_m \sim W^0$ for bounded $W$,
but the $(W-k)$-th order statistic concentrates around the quantile
$F^{-1}(1-k/W)$, which for Pareto is $(W/k)^{1/\alpha} x_m = 4^{1/\alpha} x_m$.
This is bounded, so for fixed $W$, the perturbation is $O(M/x_m)$
--- a constant.

For the perturbation to be negligible relative to the i.i.d.\ estimation
error $O(\alpha/\sqrt{k})$, we need:
\[
  \frac{\alpha^2 M}{Y_{(W-k)}} \ll \frac{\alpha}{\sqrt{k}}
  \quad\Longleftrightarrow\quad
  \frac{\alpha M \sqrt{k}}{Y_{(W-k)}} \ll 1.
\]
With $Y_{(W-k)} \sim (W/k)^{1/\alpha}$: this requires
$(W/k)^{1/\alpha} \gg \alpha M \sqrt{k}$, i.e.,
$W \gg k(\alpha M \sqrt{k})^\alpha$.
For $\alpha = 1.5$, $M = 1$, $k = 25$: requires $W \gg 25 \cdot 5^{1.5} \approx 280$.
So $W = 100$ is marginal, but clamping provides the safety net.
\end{remark}


% ===========================================================================
\section{Coupled Calibration: Tail-Index-Aware Fixed Schedule}
\label{sec:coupled}
% ===========================================================================

The implementation of $\alpha$-Robust SGD analyzed in Section~\ref{sec:adaptive}
sets $\tau_t = C \cdot \shat_t$, in which the tail-index estimate $\phat_t$
does not enter the threshold (it is purely diagnostic). The PyTorch deployment
used in our Phase~2 / Phase~3 experiments instead adopts a
\emph{tail-index-aware} fixed schedule:
\begin{equation}\label{eq:torch_fixed}
  \tau_t \;=\; \varphi(\phat_t) \cdot C \cdot \shat_t,
  \qquad \varphi(q) := q/2.
\end{equation}
The factor $\varphi(q) = q/2$ tightens the clip when noise is heavier
($p \to 1^+$, $\varphi(p) \to 1/2$) and relaxes toward the plain median
threshold under near-sub-Gaussian noise ($p \to 2$, $\varphi(p) \to 1$).
This section extends the analysis of Theorem~\ref{thm:adaptive} to the
coupled threshold~\eqref{eq:torch_fixed}, in which both $\phat_t$ and
$\shat_t$ feed into $\tau_t$.

\subsection{Calibration map}
\label{ssec:phi}

We work with a general bounded Lipschitz calibration map; the constant
choice (Section~\ref{sec:adaptive}) and the Torch
implementation~\eqref{eq:torch_fixed} are special cases.

\begin{assumption}[Calibration map]\label{ass:phi}
Let $\varphi \colon [p_{\min}, p_{\max}] \to (0, \infty)$ satisfy:
\begin{enumerate}[label=(\roman*),noitemsep]
  \item \textbf{Bounded:} there exist $0 < \varphi_{\min} \leq \varphi_{\max} < \infty$
    such that $\varphi(q) \in [\varphi_{\min}, \varphi_{\max}]$ for all
    $q \in [p_{\min}, p_{\max}]$.
  \item \textbf{Lipschitz:} there exists $L_\varphi \geq 0$ such that
    $\abs{\varphi(q_1) - \varphi(q_2)} \leq L_\varphi \abs{q_1 - q_2}$
    for all $q_1, q_2 \in [p_{\min}, p_{\max}]$.
\end{enumerate}
We further assume the true $p \in [p_{\min}, p_{\max}]$, which is automatic
when $[p_{\min}, p_{\max}] = [1.01, 1.99]$ (the implementation default) and
$p \in (1, 2)$ (Assumption~\ref{ass:heavy}).
\end{assumption}

\begin{remark}[Two specializations]\label{rem:phi_specials}
\begin{enumerate}[label=(\alph*),noitemsep,leftmargin=*]
  \item \emph{Constant calibration} $\varphi \equiv 1$: $L_\varphi = 0$,
    $\varphi_{\min} = \varphi_{\max} = 1$. Recovers $\tau_t = C\shat_t$, i.e.,
    the threshold analyzed in Theorem~\ref{thm:adaptive} (NumPy
    \texttt{optimizers/alpha\_robust\_sgd.py}).
  \item \emph{Half-index calibration} $\varphi(q) = q/2$:
    $L_\varphi = 1/2$, $\varphi_{\min} = p_{\min}/2$, $\varphi_{\max} = p_{\max}/2$.
    With $[p_{\min}, p_{\max}] = [1.01, 1.99]$, this gives
    $(\varphi_{\min}, \varphi_{\max}) = (0.505, 0.995)$. Recovers the Torch
    fixed schedule (\texttt{optimizers/torch/alpha\_robust\_sgd.py}, the
    \texttt{tail\_factor} in \texttt{\_compute\_tau\_and\_lr}).
\end{enumerate}
\end{remark}

\subsection{Joint multiplicative perturbation of \texorpdfstring{$\tau_t$}{the threshold}}
\label{ssec:joint}

Define the oracle and algorithm thresholds:
\begin{equation}\label{eq:tau_def}
  \tau_t^{*} \;:=\; \varphi(p)\cdot C\cdot \sigma_{\mathrm{med}},
  \qquad
  \tau_t \;:=\; \varphi(\phat_t)\cdot C\cdot \shat_t.
\end{equation}
The threshold ratio decomposes into independent calibration and scale errors:
\[
  \frac{\tau_t}{\tau_t^{*}}
  = \underbrace{\frac{\varphi(\phat_t)}{\varphi(p)}}_{\text{calibration}}
    \cdot \underbrace{\frac{\shat_t}{\sigma_{\mathrm{med}}}}_{\text{scale}}.
\]

\begin{lemma}[Relative calibration error]\label{lem:phi_rel}
Under Assumption~\ref{ass:phi} and for $\phat_t \in [p_{\min}, p_{\max}]$:
\begin{equation}\label{eq:rel_calib_bound}
  \abs{\frac{\varphi(\phat_t)}{\varphi(p)} - 1}
  \;\leq\; \frac{L_\varphi}{\varphi_{\min}} \cdot \abs{\phat_t - p}.
\end{equation}
For the half-index calibration $\varphi(q) = q/2$,
$L_\varphi/\varphi_{\min} = 1/p_{\min}$, so the relative error is bounded by
$\abs{\phat_t - p}/p_{\min}$.
\end{lemma}

\begin{proof}
By Lipschitz continuity (Assumption~\ref{ass:phi}(ii)),
$\abs{\varphi(\phat_t) - \varphi(p)} \leq L_\varphi\abs{\phat_t - p}$.
Dividing by $\varphi(p) \geq \varphi_{\min} > 0$ gives~\eqref{eq:rel_calib_bound}.
The half-index specialization uses $L_\varphi = 1/2$ and
$\varphi_{\min} = p_{\min}/2$.
\end{proof}

\begin{proposition}[Joint perturbation of the calibrated threshold]
\label{prop:joint_perturb}
Under Assumptions~\ref{ass:smooth}--\ref{ass:tail} and~\ref{ass:phi},
fix $\epsilon \in (0, 1/2)$ and $\Delta \in (0, \varphi_{\min}/(2 L_\varphi))$
\textup{(}so that the calibration deviation does not exceed $1/2$\textup{)}.
Define the high-probability events:
\begin{align*}
  \mathcal{E}_t &:= \bigl\{\,\abs{\shat_t/\sigma_{\mathrm{med}} - 1} \leq \epsilon\,\bigr\}
    && \text{(median concentration)}, \\
  \mathcal{H}_t &:= \bigl\{\,\abs{\phat_t - p} \leq \Delta\,\bigr\}
    && \text{(tail-index concentration)}.
\end{align*}
Set
\begin{equation}\label{eq:mu_max}
  \mu_{\max}
  \;:=\; \frac{L_\varphi}{\varphi_{\min}}\Delta \;+\; \epsilon \;+\;
    \frac{L_\varphi}{\varphi_{\min}}\Delta\cdot\epsilon.
\end{equation}
On the event $\mathcal{E}_t \cap \mathcal{H}_t$, the threshold ratio satisfies
\begin{equation}\label{eq:tau_envelope}
  \tau_t \;\in\; \bigl[(1-\mu_{\max})\tau_t^{*},\, (1+\mu_{\max})\tau_t^{*}\bigr]
\end{equation}
deterministically, equivalently $\abs{\tau_t/\tau_t^{*} - 1} \leq \mu_{\max}$.

\medskip
\textbf{Probability bookkeeping.}
By Proposition~\ref{prop:scale},
$\Prob(\mathcal{E}_t^{c}) \leq 2\exp\!\left(-2W c_{\mathrm{cdf}}^2 \epsilon^2 \sigma_{\mathrm{med}}^2\right)$.
The bound on $\Prob(\mathcal H_t^c)$ depends on the noise model:

\begin{itemize}[noitemsep,leftmargin=*]
  \item \textbf{Exact Pareto noise} \textup{(}as in Phase~1, where
    $\norm{\varepsilon_s}\sim\mathrm{Pareto}(\alpha, x_m)$ i.i.d.\
    independent of $x_s$, so $\norm{g_s}$ is a deterministic-shift of
    Pareto whose top order statistics agree with those of Pareto for large
    enough thresholds\textup{)}: Lemma~\ref{lem:hill_conc} gives
    $\Prob(\mathcal H_t^c) \leq 2\exp(-k\Delta^2/(16\alpha^2))$ for the
    Hill estimator computed on the lagged window of clean noise norms,
    plus a deterministic perturbation
    $\delta^{\mathrm{nb}}_{\mathrm{Hill}}$ from
    Proposition~\ref{prop:hill_noniid} controlling the gradient-norm vs.\
    noise-norm gap. Set
    $\Delta := \Delta_{\mathrm{conc}} + \delta^{\mathrm{nb}}_{\mathrm{Hill}}$.
  \item \textbf{General regularly-varying noise} \textup{(}beyond exact
    Pareto, e.g., the Phase~2/3 settings\textup{)}: the i.i.d.\ Pareto
    Bernstein argument of Lemma~\ref{lem:hill_conc} \emph{does not directly
    apply}; finite-sample concentration in this regime is an open question
    in the EVT literature. We therefore fall back to the deterministic
    clamping bound $\Delta = \max(p_{\max}-p, p-p_{\min})$ from
    eq.~\eqref{eq:clamp_bound}; $\mu_{\max}$ is then a non-vanishing
    constant, and Theorem~\ref{thm:coupled} below converges to a
    neighborhood of stationarity rather than to the sharp adaptive rate.
\end{itemize}

Union bound over $t = W+1, \ldots, T$:
\begin{equation}\label{eq:union_prob}
  \Prob\!\left(\bigcap_{t=W+1}^T (\mathcal{E}_t \cap \mathcal{H}_t)\right)
  \;\geq\; 1 \;-\; 2T\exp\!\left(-2W c_{\mathrm{cdf}}^2\epsilon^2\sigma_{\mathrm{med}}^2\right)
  \;-\; 2T\exp\!\left(-\tfrac{k\Delta_{\mathrm{conc}}^2}{16\alpha^2}\right),
\end{equation}
where the second term is dropped under the clamping fallback above.
\end{proposition}

\begin{proof}
On $\mathcal{E}_t \cap \mathcal{H}_t$, write the multiplicative components:
\[
  \frac{\varphi(\phat_t)}{\varphi(p)} = 1 + \kappa_t,
  \qquad
  \frac{\shat_t}{\sigma_{\mathrm{med}}} = 1 + \nu_t.
\]
By Lemma~\ref{lem:phi_rel},
$\abs{\kappa_t} \leq (L_\varphi/\varphi_{\min})\Delta$ on $\mathcal{H}_t$;
by definition, $\abs{\nu_t} \leq \epsilon$ on $\mathcal{E}_t$.
Multiplying:
\[
  \frac{\tau_t}{\tau_t^{*}}
  = (1+\kappa_t)(1+\nu_t)
  = 1 + \kappa_t + \nu_t + \kappa_t\nu_t.
\]
Therefore
\[
  \abs{\frac{\tau_t}{\tau_t^{*}} - 1}
  \leq \abs{\kappa_t} + \abs{\nu_t} + \abs{\kappa_t}\abs{\nu_t}
  \leq \frac{L_\varphi}{\varphi_{\min}}\Delta + \epsilon
    + \frac{L_\varphi}{\varphi_{\min}}\Delta\epsilon
  = \mu_{\max},
\]
which gives~\eqref{eq:tau_envelope}. The probabilistic bound
follows from Proposition~\ref{prop:scale}, Lemma~\ref{lem:hill_conc}, and a
union bound.
\end{proof}

\begin{remark}[Asymptotic order of $\mu_{\max}$]\label{rem:mu_size}
With $\epsilon = O(\sqrt{\log T/W})$ from Proposition~\ref{prop:scale} and
$\Delta = O(\alpha/\sqrt{k})$ from Lemma~\ref{lem:hill_conc} (taking
$k = W/4$, both contributions are $O(1/\sqrt{W})$ up to log factors), we obtain
\[
  \mu_{\max} = O\!\left(\sqrt{\log T / W}\right),
\]
the same asymptotic order as the median-only overhead in
Theorem~\ref{thm:adaptive}. The cross term $\kappa_t\nu_t$ is second-order
$O(1/W)$ and asymptotically negligible.

When the regularly-varying tail assumption~\ref{ass:tail} is dropped, the
deterministic clamping bound $\abs{\phat_t - p} \leq \max(p_{\max}-p, p-p_{\min})
\leq 0.99$ replaces the Hill concentration; $\mu_{\max}$ then becomes a
non-vanishing constant, and the conclusion of
Theorem~\ref{thm:coupled} below becomes convergence to a neighborhood of
stationarity rather than the sharp adaptive rate.
\end{remark}

\subsection{Coupled convergence theorem}
\label{ssec:coupled_main}

The next theorem is the central technical contribution of this section:
the coupled threshold~\eqref{eq:torch_fixed} achieves the same rate as the
oracle fixed schedule, with multiplicative constants modified by a
$1 + O(\mu_{\max})$ factor that vanishes as $W \to \infty$.

\begin{theorem}[Coupled adaptive fixed schedule]\label{thm:coupled}
Under Assumptions~\ref{ass:smooth}--\ref{ass:phi}, run $\alpha$-Robust SGD
with coupled threshold $\tau_t = \varphi(\phat_t)\cdot C\cdot \shat_t$ for
$t > W$ and a fixed $\tau_0 \geq 2M$ for $t \leq W$, with constant learning
rate $\eta$. Choose $\epsilon$ and $\Delta$ as in
Proposition~\ref{prop:joint_perturb}, with $\mu_{\max} \leq 1/2$ from~\eqref{eq:mu_max}.
Assume the regime $\tau_t^{*} \geq 4M$ for all $t > W$ (this holds, in
particular, in the optimal $\tau^{*} = \Theta(T^{1/(3p-2)})$ regime once $T$
exceeds a problem-dependent constant).

Define the calibrated bias and variance constants:
\[
  B^{*}_\varphi := B(\varphi(p)\cdot C\cdot \sigma_{\mathrm{med}}),
  \qquad
  V^{*}_\varphi := V(\varphi(p)\cdot C\cdot \sigma_{\mathrm{med}}),
\]
where $B(\cdot)$ and $V(\cdot)$ are from Definition~\ref{def:BV}.

On the high-probability event $\mathcal{G} := \bigcap_{t=W}^T (\mathcal{E}_t \cap \mathcal{H}_t)$
\textup{(}whose complement is bounded by~\eqref{eq:union_prob}\textup{)}:
\begin{equation}\label{eq:coupled_bound}
  \frac{1}{T}\sum_{t=1}^T \E\!\left[\norm{\nabla f(x_t)}^2 \;\Big|\; \mathcal{G}\right]
  \;\leq\;
    \frac{2\Delta_f}{\eta T}
    + (1 + \delta_W^\varphi)^2\,(B^{*}_\varphi)^2
    + (1 + \delta_W^\varphi)\, L\eta\,(V^{*}_\varphi)^2
    + O\!\left(\frac{W}{T}\right),
\end{equation}
where the calibration overhead is
\begin{equation}\label{eq:delta_W_phi}
  \delta_W^\varphi
  := \max\!\Bigl(4(p-1)\mu_{\max},\; 2(2-p)\mu_{\max}\Bigr)
  \;\leq\; 4\mu_{\max}
  \;=\; O\!\Bigl(\sqrt{\log T/W}\Bigr).
\end{equation}

\medskip
\textbf{Optimal scaling.}
Choosing $C$ and $\eta$ to balance the bias and variance terms (as in
Theorem~\ref{thm:fixed} Steps~6--7) yields, on $\mathcal{G}$:
\begin{equation}\label{eq:coupled_rate}
  \frac{1}{T}\sum_{t=1}^T \E\!\left[\norm{\nabla f(x_t)}^2 \;\Big|\; \mathcal{G}\right]
  \;\leq\;
    C_{\mathrm{rate}}^{\varphi}\,(1 + O(\delta_W^\varphi))
    \cdot T^{-\frac{2(p-1)}{3p-2}}
    \;+\; O\!\left(\frac{W}{T}\right),
\end{equation}
where $C_{\mathrm{rate}}^{\varphi}$ depends on $L, \Delta_f, \sigma, M, p, \varphi(p)$
but not on $T$ or $W$. The rate exponent $-2(p-1)/(3p-2)$ is identical to
that of the oracle fixed schedule (Theorem~\ref{thm:fixed}); the calibration
map $\varphi$ enters only through multiplicative constants and a vanishing
$O(\sqrt{\log T/W})$ overhead.
\end{theorem}

\begin{proof}
\textbf{Step 1: Burn-in ($t \leq W$).}
For $t \leq W$, the threshold $\tau_0 \geq 2M$ is used. The same descent
argument as in Theorem~\ref{thm:adaptive} (burn-in step) gives total cost
$O(W)$, contributing $O(W/T)$ to~\eqref{eq:coupled_bound}.

\medskip
\textbf{Step 2: Pointwise threshold envelope on $\mathcal{G}$.}
For $t > W$, condition on $\mathcal{G}$. By
Proposition~\ref{prop:joint_perturb},
$\tau_t \in [(1-\mu_{\max})\tau_t^{*},\, (1+\mu_{\max})\tau_t^{*}]$
deterministically. Since $\tau_t^{*} \geq 4M$ and $\mu_{\max} \leq 1/2$:
$\tau_t \geq (1-\mu_{\max})\tau_t^{*} \geq 2M$, so the simplified bias
bound $B(\tau) \leq 2^{p-1}\sigma^p/\tau^{p-1}$ from
Lemma~\ref{lem:bias} applies.

\medskip
\textbf{Step 3: Inflated bias bound.}
Since $B(\tau)$ is decreasing in $\tau$, use the lower envelope
$\tau_t \geq (1-\mu_{\max})\tau_t^{*}$:
\[
  B(\tau_t) \leq \frac{2^{p-1}\sigma^p}{\tau_t^{p-1}}
  = \left(\frac{\tau_t^{*}}{\tau_t}\right)^{p-1} B(\tau_t^{*})
  \leq (1-\mu_{\max})^{-(p-1)}\, B(\tau_t^{*}).
\]
Bound the prefactor using two elementary inequalities valid for
$x \in [0, 1/2]$ and $p \in (1, 2]$:
\begin{align*}
  \log\!\frac{1}{1-x} &\leq 2x \quad\text{(since $g(x) = \log(1/(1-x)) - 2x$ is decreasing on $[0, 1/2]$ with $g(0) = 0$)}, \\
  e^y &\leq 1 + 2y \quad\text{for } y \in [0, 1] \quad\text{(convexity of $e^y$ on $[0,1]$)}.
\end{align*}
Therefore, for $\mu_{\max} \leq 1/2$:
\begin{align*}
  (1-\mu_{\max})^{-(p-1)}
  &= \exp\!\bigl((p-1)\log(1/(1-\mu_{\max}))\bigr)
  \leq \exp\!\bigl(2(p-1)\mu_{\max}\bigr)
  \leq 1 + 4(p-1)\mu_{\max},
\end{align*}
where the last step uses $2(p-1)\mu_{\max} \leq 2 \cdot 1 \cdot 1/2 = 1$.
This yields:
\begin{equation}\label{eq:bias_phi_inflate}
  B(\tau_t) \;\leq\; \bigl(1 + 4(p-1)\mu_{\max}\bigr)\, B(\tau_t^{*}).
\end{equation}

\medskip
\textbf{Step 4: Inflated variance bound.}
$V(\tau)^2$ is increasing in $\tau$. Use the upper envelope
$\tau_t \leq (1+\mu_{\max})\tau_t^{*}$:
\[
  V(\tau_t)^2 = 2^{p-1}(M^p + \sigma^p)\,\tau_t^{2-p}
  \leq (1+\mu_{\max})^{2-p}\, V(\tau_t^{*})^2.
\]
For $x \in [0, 1/2]$ and $p \in (1, 2]$ (so $2-p \in [0, 1)$):
\[
  (1+x)^{2-p} = \exp\!\bigl((2-p)\log(1+x)\bigr)
  \leq \exp\!\bigl((2-p) x\bigr)
  \leq 1 + 2(2-p)x,
\]
using $\log(1+x) \leq x$ and the same $e^y \leq 1+2y$ bound on
$y = (2-p)x \leq 1/2$. Therefore:
\begin{equation}\label{eq:var_phi_inflate}
  V(\tau_t)^2 \;\leq\; \bigl(1 + 2(2-p)\mu_{\max}\bigr)\, V(\tau_t^{*})^2.
\end{equation}

\medskip
\textbf{Step 5: One-step descent on $\mathcal{G}$.}
Plug the inflated bounds~\eqref{eq:bias_phi_inflate}--\eqref{eq:var_phi_inflate}
into the descent inequality~\eqref{eq:descent_final}, conditioning on
$\Filt_t \cap \mathcal{G}$. By the lagged-window contract
\textup{(}Remark~\ref{rem:lag}\textup{)} both $\shat_t$ and $\phat_t$ are
computed from $\{\norm{g_s}\}_{s < t}$ and are therefore
$\Filt_{t}$-measurable; the events $\mathcal{E}_t$ and $\mathcal{H}_t$ for
$t > W$ are $\Filt_t$-measurable as well. Consequently
$\mathcal G = \bigcap_{s = W+1}^T (\mathcal E_s \cap \mathcal H_s)$ is
$\Filt_T$-measurable, so $\Prob(\,\cdot\,\mid\mathcal G)$ admits a regular
conditional probability and the tower law applies as in Theorem~\ref{thm:adaptive}, Step~6.
\begin{equation}\label{eq:coupled_one_step}
  \E\!\left[f(x_{t+1}) \,\big|\, \Filt_t, \mathcal{G}\right]
  \leq f(x_t) - \frac{\eta}{2}\norm{\nabla f(x_t)}^2
    + \frac{\eta}{2}(1 + \delta_W^\varphi)^2 (B^{*}_\varphi)^2
    + \frac{L\eta^2}{2}(1 + \delta_W^\varphi)(V^{*}_\varphi)^2,
\end{equation}
where we have used $\tau_t^{*} = \varphi(p) C \sigma_{\mathrm{med}}$
(time-independent in the fixed schedule) and absorbed $4(p-1)$ vs.\ $2(2-p)$
into the unified factor $\delta_W^\varphi$ defined in~\eqref{eq:delta_W_phi}.

\medskip
\textbf{Step 6: Telescope.}
Take full expectations on $\mathcal{G}$ and sum
$t = W+1, \ldots, T$:
\[
  \E\!\left[f(x_{T+1}) \mid \mathcal{G}\right] - \E\!\left[f(x_{W+1}) \mid \mathcal{G}\right]
  \leq -\frac{\eta}{2}\sum_{t=W+1}^T \E\!\left[\norm{\nabla f(x_t)}^2 \mid \mathcal{G}\right]
    + \frac{(T-W)\eta}{2}(1 + \delta_W^\varphi)^2 (B^{*}_\varphi)^2
\]
\[
    + \frac{(T-W)L\eta^2}{2}(1 + \delta_W^\varphi)(V^{*}_\varphi)^2.
\]
Adding the burn-in cost (Step~1) and using $f(x_{T+1}) \geq f^{*}$, then
dividing by $T\eta/2$ gives~\eqref{eq:coupled_bound}.

\medskip
\textbf{Step 7: Optimization.}
Inequality~\eqref{eq:coupled_bound} differs from the oracle fixed-schedule
bound~\eqref{eq:general_bound} only by the multiplicative factors
$(1 + \delta_W^\varphi)^2$ on the bias term and $(1 + \delta_W^\varphi)$ on
the variance term, plus the additive $O(W/T)$ burn-in. The same balancing
argument as Theorem~\ref{thm:fixed} (Steps~6--7) optimizes $C$ and $\eta$ to
yield the rate $T^{-2(p-1)/(3p-2)}$, with the constant
$C_{\mathrm{rate}}^\varphi$ inheriting a factor at most
$(1 + \delta_W^\varphi)^{O(1)} = 1 + O(\delta_W^\varphi)$. The
result~\eqref{eq:coupled_rate} follows.

\medskip
\textbf{Step 8: Verifying $\tau_t^{*} \geq 4M$.}
With optimal $C^{*}$ and $\tau^{*} = \varphi(p)\, C^{*}\, \sigma_{\mathrm{med}}
= \Theta(T^{1/(3p-2)})$, the regime $\tau_t^{*} \geq 4M$ holds for all
$T \geq T_0$ with $T_0 = (4M / \varphi_{\min})^{3p-2} \cdot
\Theta(1)$, which is finite and independent of $W$. For $T < T_0$ we
fall back on the burn-in fixed threshold $\tau_0$, incurring $O(T_0) = O(1)$
additive cost folded into $O(W/T)$.
\end{proof}

\begin{corollary}[Half-index calibration]\label{cor:half_index}
Specializing Theorem~\ref{thm:coupled} to $\varphi(q) = q/2$ (the Torch
implementation~\eqref{eq:torch_fixed}) yields the convergence rate
\[
  T^{-\frac{2(p-1)}{3p-2}},
\]
identical to that of Theorem~\ref{thm:adaptive} (constant $\varphi$). The
constants are modified by powers of $\varphi(p) = p/2$:
\begin{equation}\label{eq:half_index_constants}
  (B^{*}_\varphi)^2 = \bigl(p/2\bigr)^{-2(p-1)}\, B(C\sigma_{\mathrm{med}})^2,
  \qquad
  (V^{*}_\varphi)^2 = \bigl(p/2\bigr)^{2-p}\, V(C\sigma_{\mathrm{med}})^2.
\end{equation}
For $p \in [1.01, 1.99]$, $\varphi(p) = p/2 \in [0.505, 0.995]$, so the
bias term is multiplied by at most $(0.505)^{-2(p-1)}$ and the variance
term by $(0.505)^{2-p}$; both factors are bounded constants absorbed into
$C_{\mathrm{rate}}^\varphi$.
\end{corollary}

\begin{remark}[Bias--variance trade-off interpretation]\label{rem:bv_tradeoff}
The half-index calibration $\varphi(q) = q/2$ \emph{decreases} the variance
constant by a factor of $(p/2)^{2-p}$ (since $p/2 < 1$ and $2-p \geq 0$)
and \emph{increases} the bias constant by a factor of $(p/2)^{-2(p-1)}$.
For $p = 1.5$:
\begin{itemize}[noitemsep,leftmargin=*]
  \item Variance multiplier: $(0.75)^{0.5} \approx 0.866$ (13.4\% reduction).
  \item Bias multiplier: $(0.75)^{-1.0} \approx 1.333$ (33.3\% increase).
\end{itemize}
This is consistent with the design intuition behind the
$\varphi(q) = q/2$ heuristic: under heavier tails (smaller $p$), the
calibration tightens the clip, reducing the variance contribution at the
cost of a moderate increase in bias. Whether this trade-off is favorable
depends on the relative magnitudes of $L\Delta_f$, $\sigma$, and $M$ at
the specific problem instance; in the asymptotic rate analysis, both
contributions are absorbed into $C_{\mathrm{rate}}^\varphi$.
\end{remark}

\begin{remark}[Theory--code alignment]\label{rem:code_align}
Theorem~\ref{thm:coupled} closes the analytical gap between the implementation
variants of $\alpha$-Robust SGD documented in
\texttt{submission/proof\_code\_matrix.md}:
\begin{itemize}[noitemsep,leftmargin=*]
  \item \textbf{NumPy Phase~1} (\texttt{optimizers/alpha\_robust\_sgd.py},
    \texttt{schedule="fixed"}): $\varphi \equiv 1$. Recovered as the
    $L_\varphi = 0$ specialization of Theorem~\ref{thm:coupled}, equivalent
    to Theorem~\ref{thm:adaptive}.
  \item \textbf{Torch Phase~2/3} (\texttt{optimizers/torch/alpha\_robust\_sgd.py},
    \texttt{schedule="fixed"}, \texttt{tail\_factor = p\_hat\_val/2}):
    $\varphi(q) = q/2$. Covered by Corollary~\ref{cor:half_index}.
  \item \textbf{Future variants:} any bounded Lipschitz $\varphi$ on
    $[p_{\min}, p_{\max}]$ inherits the same rate exponent. Implementers may
    tune $\varphi$ for empirical performance without invalidating the analysis.
\end{itemize}
\end{remark}

\begin{remark}[Why the rate is invariant to $\varphi$]\label{rem:phi_rate}
The rate exponent $-2(p-1)/(3p-2)$ in Theorem~\ref{thm:fixed} arises from
balancing two power-law functions of $\tau$: the bias scales as
$\tau^{-(p-1)}$ and the variance as $\tau^{2-p}$. A multiplicative
calibration $\tau \mapsto \varphi(p)\tau$ shifts both functions by constant
factors but preserves the \emph{exponents} of the power laws, hence the
optimal balance point $\tau^{*}$ is shifted by a constant factor and the
rate exponent is unchanged. The calibration map therefore enters as a
constant-factor calibration, absorbed into $C_{\mathrm{rate}}^\varphi$;
it cannot improve or worsen the asymptotic rate.

This invariance is the technical content of the \emph{coupled} extension:
even though both $\phat_t$ and $\shat_t$ feed into $\tau_t$, the
estimator-induced perturbations enter the bound only through the bounded
multiplicative envelope~\eqref{eq:tau_envelope} and the inflation factors
of size $1 + O(\mu_{\max})$.
\end{remark}


% ===========================================================================
\section{Growing-Schedule Convergence}
\label{sec:growing}
% ===========================================================================

Take $\tau_t = \shat_t \cdot t^{1/\phat_t}$ and $\eta_t = \eta_0/\sqrt{t}$ (growing schedule).

\begin{lemma}[Threshold drift decomposition]\label{lem:drift}
Let $\tau_t^* = \sigma \cdot t^{1/p}$ be the oracle threshold and
$\tau_t = \shat_t \cdot t^{1/\phat_t}$ the $\alpha$-Robust threshold.
Define the multiplicative drift factor $D_t := \tau_t / \tau_t^*$ and the
estimation error $\Delta_t := \phat_t - p$. Then:
\begin{equation}
  D_t = \frac{\shat_t}{\sigma} \cdot t^{1/\phat_t - 1/p}
  = \frac{\shat_t}{\sigma} \cdot t^{-\Delta_t/(\phat_t \cdot p)}.
\end{equation}
\end{lemma}

\begin{proof}
Direct computation:
\[
  \frac{1}{\phat_t} - \frac{1}{p}
  = \frac{p - \phat_t}{p \cdot \phat_t}
  = \frac{-\Delta_t}{p \cdot \phat_t}.
\]
Therefore $t^{1/\phat_t - 1/p} = t^{-\Delta_t/(p \cdot \phat_t)}$.
\end{proof}

\begin{lemma}[Drift factor bound]\label{lem:drift_bound}
For $\abs{\Delta_t} \leq \Delta_{\max}$ (guaranteed by the clamping~\eqref{eq:clamp_bound})
and $\shat_t/\sigma \in [c_1, c_2]$ (from Proposition~\ref{prop:scale}):
\[
  \abs{\log D_t} \leq \abs{\log(c_2)} + \frac{\Delta_{\max} \log t}{p_{\min}^2},
\]
so $D_t \in [D_{\min}, D_{\max}]$ with:
\begin{equation}\label{eq:Dmax}
  D_{\max} = c_2 \cdot T^{\Delta_{\max}/p_{\min}^2}.
\end{equation}
With Hill estimator concentration (Lemma~\ref{lem:hill_conc}),
$\Delta_{\max} = O(\alpha/\sqrt{k})$ with high probability, giving
$D_{\max} = T^{O(1/\sqrt{W})}$.
\end{lemma}

\begin{proof}
Taking logarithms of $D_t = (\shat_t/\sigma) \cdot t^{-\Delta_t/(p\phat_t)}$:
\[
  \log D_t = \log(\shat_t/\sigma) - \frac{\Delta_t}{p\phat_t}\log t.
\]
Bounding each term: $\abs{\log(\shat_t/\sigma)} \leq \abs{\log c_2}$
(or $\abs{\log c_1}$, whichever is larger), and
$\abs{\Delta_t/(p\phat_t)} \leq \Delta_{\max}/p_{\min}^2$. The result follows.
\end{proof}

\begin{theorem}[{Growing-schedule convergence --- all $p \in (1, 2]$}]\label{thm:growing}
Under Assumptions~\ref{ass:smooth}--\ref{ass:tail}, let $\alpha$-Robust SGD
run for $T$ steps with $\eta_t = \eta_0/\sqrt{t}$ and
$\tau_t = \shat_t \cdot t^{1/\phat_t}$ (lagged window;
Remark~\ref{rem:lag}). Adopt the following standing conventions:
\begin{enumerate}[label=\textup{(}\roman*\textup{)},noitemsep,leftmargin=*]
  \item the oracle threshold is taken as
    $\tau_t^{*} := \sigma_{\mathrm{med}} \cdot t^{1/p}$ (using the median
    scale $\sigma_{\mathrm{med}}$, since $\shat_t \to \sigma_{\mathrm{med}}$
    by Proposition~\ref{prop:scale}; the additional moment vs.\ median
    constant $\sigma_{\mathrm{med}}/\sigma$ is a problem-dependent
    multiplier that we absorb into $C'$),
  \item the drift factor satisfies $D_t \in [1/D_{\max}, D_{\max}]$ for all
    $t > W$ (which holds with high probability by Lemma~\ref{lem:drift_bound};
    deterministically with $\Delta_{\max} = 0.99$ from clamping).
\end{enumerate}

\medskip
Then for \textbf{every} $p \in (1, 2]$ and for the
$T$-dependent choice $\eta_0 = \eta_0^{*}(T)$ defined in Step~6 below:
\begin{equation}\label{eq:growing_main}
  \min_{1 \leq t \leq T}\E\!\left[\norm{\nabla f(x_t)}^2\right]
  \leq C'(D_{\max}, p, L, \sigma_{\mathrm{med}}, M, \Delta_f) \cdot T^{-\frac{p-1}{p}},
\end{equation}
where $C'$ depends polynomially on $D_{\max}$ and on the problem constants,
but the rate exponent $-(p-1)/p$ is universal over $p \in (1, 2]$.

\medskip
\textbf{Known-horizon caveat.} The optimal $\eta_0^{*} \propto T^{(p-2)/(2p)}$
depends on the horizon~$T$ (Step~6 of the proof). The growing-threshold
schedule \emph{cancels} the $T$-dependence in $\tau_t$, but \emph{not} in
the learning rate. Genuine ``unknown-$T$'' deployment requires either
\textup{(a)} the doubling trick (Remark~\ref{rem:horizon}), incurring a
$\log T$ overhead, or \textup{(b)} replacing $\eta_t = \eta_0/\sqrt t$ with
$\eta_t = \eta_0/t^{p/(2p-2)}$, which yields the rate without horizon
knowledge at the cost of a slightly different constant.
\end{theorem}

\begin{proof}
\textbf{Step 1: Descent with time-varying parameters.}
At each step $t > W$, by the same argument as Theorem~\ref{thm:fixed} Steps 1--4:
\[
  \E[f(x_{t+1}) \mid \Filt_t]
  \leq f(x_t) - \frac{\eta_t}{2}\norm{\nabla f(x_t)}^2
    + \frac{\eta_t}{2}B(\tau_t)^2
    + \frac{L\eta_t^2}{2}V(\tau_t)^2.
\]

\textbf{Step 2: Perturbed bounds via drift factor.}
Since $\tau_t = D_t \cdot \tau_t^*$:
\[
  B(\tau_t) = \frac{2^{p-1}\sigma^p}{(\tau_t - M)^{p-1}}
  \leq \frac{2^{p-1}\sigma^p}{(D_{\min}\tau_t^* - M)^{p-1}}.
\]
For $D_{\min}\tau_t^* \geq 2M$ (which holds for $t \geq t_0$ since
$\tau_t^* = \sigma t^{1/p} \to \infty$):
\[
  B(\tau_t) \leq \frac{2^{2(p-1)}\sigma^p}{(D_{\min}\tau_t^*)^{p-1}}
  = D_{\max}^{p-1} \cdot B(\tau_t^*).
\]
Similarly: $V(\tau_t)^2 \leq D_{\max}^{2-p} \cdot V(\tau_t^*)^2$.

Define the \emph{drift penalty}:
$\Gamma_B := D_{\max}^{2(p-1)}$,\; $\Gamma_V := D_{\max}^{2-p}$.

\textbf{Step 3: Oracle bias and variance.}
With $\tau_t^* = \sigma t^{1/p}$:
\[
  B(\tau_t^*)^2 = \frac{A_B}{\sigma^{2(p-1)}} \cdot t^{-2(p-1)/p}
  =: C_B \cdot t^{-2(p-1)/p},
\]
\[
  V(\tau_t^*)^2 = A_V \sigma^{2-p} \cdot t^{(2-p)/p}
  =: C_V \cdot t^{(2-p)/p},
\]
where $C_B = A_B/\sigma^{2(p-1)}$ and $C_V = A_V\sigma^{2-p}$ are
distribution-dependent constants.

\textbf{Step 4: Telescope.}
Taking full expectations and summing for $t = t_0, \ldots, T$:
\begin{equation}\label{eq:growing_tele}
  \sum_{t=t_0}^T \frac{\eta_t}{2}\E[\norm{\nabla f(x_t)}^2]
  \leq \Delta_f
    + \Gamma_B \sum_{t=t_0}^T \frac{\eta_t}{2} C_B t^{-2(p-1)/p}
    + \Gamma_V \sum_{t=t_0}^T \frac{L\eta_t^2}{2} C_V t^{(2-p)/p}.
\end{equation}

\textbf{Step 5: Evaluate all sums (unified for all $p \in (1,2]$).}

\emph{Left side.}
By the integral comparison $\sum_{t=1}^T t^{-1/2} \geq \int_1^{T+1} t^{-1/2}\,dt
= 2\sqrt{T+1} - 2$:
\begin{equation}\label{eq:left_sum}
  \sum_{t=1}^T \frac{\eta_0}{2\sqrt{t}} \geq \eta_0(\sqrt{T} - 1).
\end{equation}

\emph{Sum~S$_1$ (bias contribution).}
\[
  S_1 := \sum_{t=1}^T \frac{\eta_0}{2\sqrt{t}} \cdot C_B\, t^{-2(p-1)/p}
  = \frac{\eta_0 C_B}{2} \sum_{t=1}^T t^{-\beta_1},
\]
where $\beta_1 := \frac{1}{2} + \frac{2(p-1)}{p} = \frac{p + 4p - 4}{2p}
= \frac{5p - 4}{2p}$.

By the integral comparison $\sum_{t=1}^T t^{-\beta}
\leq 1 + \int_1^T t^{-\beta}\,dt$:
\begin{equation}\label{eq:S1_all_p}
  \sum_{t=1}^T t^{-\beta_1} \leq
  \begin{cases}
    \frac{T^{1-\beta_1}}{1-\beta_1} + 1
      & \text{if } \beta_1 < 1 \;(p < 4/3), \\[4pt]
    \ln T + 1
      & \text{if } \beta_1 = 1 \;(p = 4/3), \\[4pt]
    \frac{1}{\beta_1 - 1} + 1 =: \zeta_1(p)
      & \text{if } \beta_1 > 1 \;(p > 4/3).
  \end{cases}
\end{equation}

\emph{Sum~S$_2$ (variance contribution).}
\[
  S_2 := \sum_{t=1}^T \frac{L\eta_0^2}{2t} \cdot C_V\, t^{(2-p)/p}
  = \frac{L\eta_0^2 C_V}{2} \sum_{t=1}^T t^{-\beta_2},
\]
where $\beta_2 := 1 - \frac{2-p}{p} = \frac{2(p-1)}{p}$.
For all $p \in (1, 2)$: $\beta_2 \in (0, 1)$, so:
\begin{equation}\label{eq:S2_bound}
  \sum_{t=1}^T t^{-\beta_2} \leq \frac{T^{1-\beta_2}}{1-\beta_2} + 1
  = \frac{p \cdot T^{(2-p)/p}}{2-p} + 1.
\end{equation}
For $p = 2$: $\beta_2 = 1$, so $\sum t^{-1} \leq \ln T + 1$.

\textbf{Step 6: Extract convergence rate (unified).}
Dividing~\eqref{eq:growing_tele} by $\eta_0(\sqrt{T}-1) \geq \eta_0\sqrt{T}/2$
(for $T \geq 4$):
\begin{equation}\label{eq:growing_three_terms}
  \min_t \E[\norm{\nabla f}^2] \leq
  \underbrace{\frac{2\Delta_f}{\eta_0\sqrt{T}}}_{\text{(I)}}
  + \underbrace{\Gamma_B \cdot \frac{C_B \cdot \Sigma_1(T)}{\sqrt{T}}}_{\text{(II)}}
  + \underbrace{\Gamma_V \cdot \frac{LC_V\eta_0 \cdot \Sigma_2(T)}{\sqrt{T}}}_{\text{(III)}},
\end{equation}
where $\Sigma_j(T) := \sum_{t=1}^T t^{-\beta_j}$.

\emph{Term~(II): bias contribution.}
Using the bounds~\eqref{eq:S1_all_p}:
\begin{itemize}[noitemsep]
  \item For $p > 4/3$: $\Sigma_1 = O(1)$, so (II) $= O(T^{-1/2})$.
  \item For $p = 4/3$: $\Sigma_1 = O(\ln T)$, so (II) $= O(\ln T / \sqrt{T})$.
  \item For $1 < p < 4/3$: $\Sigma_1 = O(T^{(4-3p)/(2p)})$, so
    (II) $= O(T^{(4-3p)/(2p) - 1/2}) = O(T^{-2(p-1)/p})$.
\end{itemize}

In each regime, (II) $= O(T^{-2(p-1)/p})$ or smaller.
Since $2(p-1)/p > (p-1)/p$ for $p > 1$, term~(II) is always dominated by the
final rate $T^{-(p-1)/p}$ (established below). Therefore term~(II) is
asymptotically negligible.

\emph{Terms~(I) and~(III): balancing.}
Using $\Sigma_2(T) = O(T^{(2-p)/p})$ (or $O(\ln T)$ for $p = 2$),
terms (I) and (III) become:
\[
  \text{(I)} = \frac{2\Delta_f}{\eta_0\sqrt{T}}, \qquad
  \text{(III)} = \Gamma_V \frac{LC_V\eta_0 p}{2-p} \cdot T^{(2-p)/p - 1/2}
  = \Gamma_V \frac{LC_V\eta_0 p}{2-p} \cdot T^{(4-3p)/(2p)}.
\]
Setting (I) $=$ (III):
\[
  \frac{2\Delta_f}{\eta_0\sqrt{T}}
  = \Gamma_V \frac{LC_V\eta_0 p}{2-p} T^{(4-3p)/(2p)}
  \;\Longrightarrow\;
  \eta_0^2 = \frac{2\Delta_f(2-p)}{\Gamma_V LC_V p}
    \cdot T^{-1/2 - (4-3p)/(2p)}.
\]
The exponent on $T$ is:
\[
  -\frac{1}{2} - \frac{4-3p}{2p} = -\frac{p + 4 - 3p}{2p}
  = -\frac{4-2p}{2p} = \frac{p-2}{p}.
\]
So $\eta_0^* = \sqrt{\frac{2\Delta_f(2-p)}{\Gamma_V LC_V p}} \cdot T^{(p-2)/(2p)}$
(which decreases for $p < 2$, ensuring $\eta_0 \to 0$).

Substituting back into (I):
\[
  \text{(I)} = \frac{2\Delta_f}{\eta_0^*\sqrt{T}}
  = \sqrt{\frac{2\Delta_f \Gamma_V LC_V p}{2-p}}
    \cdot T^{(2-p)/(2p) - 1/2}
  = \sqrt{\frac{2\Delta_f \Gamma_V LC_V p}{2-p}}
    \cdot T^{-(p-1)/p}.
\]
This confirms the rate is $T^{-(p-1)/p}$.

\textbf{Step 7: Verify term~(II) is dominated.}
We claimed (II) is dominated by the $T^{-(p-1)/p}$ rate. Check:
$2(p-1)/p > (p-1)/p$ iff $p-1 > 0$ iff $p > 1$. \checkmark

For $p > 4/3$: (II) $= O(T^{-1/2})$. Compare with $(p-1)/p$:
$1/2 > (p-1)/p$ iff $p > 2p-2$ iff $p < 2$. \checkmark

So for all $p \in (1, 2)$, term~(II) decays strictly faster than the rate.
At $p = 2$: both are $T^{-1/2}$ (up to $\log$ factors), consistent.

\textbf{Step 8: Compare with fixed schedule.}
The growing-schedule rate $T^{-(p-1)/p}$ vs.\ the fixed-schedule rate
$T^{-2(p-1)/(3p-2)}$:
\[
  \frac{p-1}{p} < \frac{2(p-1)}{3p-2}
  \;\;\Longleftrightarrow\;\;
  (p-1)(3p-2) < 2p(p-1)
  \;\;\Longleftrightarrow\;\;
  3p - 2 < 2p
  \;\;\Longleftrightarrow\;\;
  p < 2.
\]
So for all $p \in (1, 2)$, the fixed schedule achieves a \emph{strictly better}
rate. At $p = 2$, both rates equal $T^{-1/2}$.

\begin{center}
\renewcommand{\arraystretch}{1.2}
\begin{tabular}{cccc}
\toprule
$p$ & Fixed rate & Growing rate & Gap \\
\midrule
$1.2$ & $T^{-1/4}$ & $T^{-1/6}$ & $T^{-1/12}$ \\
$1.5$ & $T^{-2/5}$ & $T^{-1/3}$ & $T^{-1/15}$ \\
$1.8$ & $T^{-8/17}$ & $T^{-4/9}$ & $T^{-4/153}$ \\
$2.0$ & $T^{-1/2}$ & $T^{-1/2}$ & none \\
\bottomrule
\end{tabular}
\end{center}

\textbf{Step 9: Role of $D_{\max}$.}
By Lemma~\ref{lem:drift_bound}, $D_{\max} = T^{O(\Delta_{\max}/p_{\min}^2)}$.
\begin{itemize}[noitemsep]
  \item With clamping only: $\Delta_{\max} \leq 0.99$, so
    $D_{\max} \leq T^{0.99/p_{\min}^2}$. The drift penalties
    $\Gamma_B = D_{\max}^{2(p-1)}$ and $\Gamma_V = D_{\max}^{2-p}$ are
    polynomial in $T$ with exponents at most $2 \cdot 0.99/p_{\min}^2 \approx 1.94$.
    These enter as multiplicative constants that do not affect the rate exponent.
  \item With Hill concentration (Lemma~\ref{lem:hill_conc}):
    $\Delta_{\max} = O(4\alpha/\sqrt{k})$ with probability
    $\geq 1 - 2\exp(-k/(16))$. For $k = W/4$: $D_{\max} = T^{O(1/\sqrt{W})}$.
\end{itemize}
\end{proof}

\begin{corollary}[Sufficient window size]\label{cor:window}
For the growing schedule, the drift penalties affect the constant $C'$
but not the rate. To ensure $C' \leq C_0 \cdot T^{\gamma}$ with
$\gamma < (p-1)/(2p)$ (so the drift penalty is dominated by the rate),
we need $\Gamma_V = D_{\max}^{2-p} \leq T^{(p-1)/(2p)}$, i.e.,
$(2-p)\Delta_{\max}/p_{\min}^2 < (p-1)/(2p)$.

\emph{With clamping} ($\Delta_{\max} = 0.99$): this requires
$(2-p) \cdot 0.99 / p_{\min}^2 < (p-1)/(2p)$. For $p = 1.5$:
$0.5 \cdot 0.99/1.02 < 1/(6) \Rightarrow 0.486 < 0.167$ --- \textbf{fails}.
So clamping alone gives a large constant for $p < 1.8$.

\emph{With Hill estimation} ($\Delta_{\max} = c\alpha/\sqrt{k}$):
the constraint becomes $W \geq (2p(2-p)c\alpha / ((p-1)p_{\min}^2))^2$.

\begin{center}
\renewcommand{\arraystretch}{1.2}
\begin{tabular}{cccc}
\toprule
$p$ & Rate exponent & Sufficient $W$ (approx.) & $W = 100$? \\
\midrule
$1.2$ & $-1/6$ & $\sim 500$ & marginal \\
$1.5$ & $-1/3$ & $\sim 150$ & close \\
$1.8$ & $-4/9$ & $\sim 50$ & sufficient \\
$2.0$ & $-1/2$ & $\sim 20$ & sufficient \\
\bottomrule
\end{tabular}
\end{center}
\end{corollary}


% ===========================================================================
\section{Discussion}
\label{sec:discussion}
% ===========================================================================

\subsection{Why $\alpha$-Robust Beats NormalizedSGD}

NormalizedSGD uses $x_{t+1} = x_t - \eta_t g_t/\norm{g_t}$, discarding gradient
magnitude entirely. While it achieves the same minimax-optimal rate in theory
(Cutkosky \& Mehta, 2021), it pays a hidden constant cost.

\textbf{Gradient magnitude carries curvature information.} In steep regions,
$\norm{\nabla f(x_t)}$ is large and larger steps are warranted. NormalizedSGD
is blind to this: it takes unit-norm steps everywhere.

\textbf{Clipping preserves magnitude for moderate gradients.} When
$\norm{g_t} \leq \tau$, $\clip(g_t, \tau) = g_t$ --- the full gradient is used.
Only extreme outliers are attenuated. With $\tau = 5\shat$, clipping occurs
on roughly $5$--$15\%$ of steps (by the heavy-tail property), preserving the
curvature signal for the remaining steps.

Let $\rho = \Prob(\norm{g_t} > \tau)$ be the clipping frequency. The effective
signal per step for $\alpha$-Robust is
$\approx (1-\rho)\norm{\nabla f(x_t)} + \rho\tau$, while NormalizedSGD's signal
is always $1$. When $\norm{\nabla f(x_t)} > 1$, $\alpha$-Robust's per-step
progress exceeds NormalizedSGD's.


\subsection{Summary of Results}

\begin{table}[h]
\centering
\renewcommand{\arraystretch}{1.3}
\begin{tabular}{lll}
\toprule
\textbf{Result} & \textbf{Rate} & \textbf{Requirements} \\
\midrule
Theorem~\ref{thm:fixed} (fixed, oracle $\tau$)
  & $O(T^{-2(p-1)/(3p-2)})$
  & Known $T$, $p$, $\sigma$ \\
Theorem~\ref{thm:adaptive} (fixed, adaptive $\tau = C\shat$)
  & $O(T^{-2(p-1)/(3p-2)}) + O(\sqrt{\log T/W})$
  & Known $T$; $\shat_t$ estimates $\sigma$ \\
Theorem~\ref{thm:coupled} (fixed, coupled $\tau = \varphi(\phat) C\shat$)
  & $O(T^{-2(p-1)/(3p-2)}) + O(\sqrt{\log T/W})$
  & Known $T$; bounded Lipschitz $\varphi$;
    $\shat_t,\phat_t$ joint estimation \\
Theorem~\ref{thm:growing} (growing $\tau$)
  & $O(T^{-(p-1)/p})$
  & See theorem for $\eta_0(T)$; clamping if Hill off-model \\
\bottomrule
\end{tabular}
\caption{Summary of convergence guarantees.}
\end{table}

\begin{enumerate}
  \item Fixed $\tau \propto T^{1/(3p-2)}$ hits $T^{-2(p-1)/(3p-2)}$
    \textup{(}Zhang et al., 2020\textup{)}.
  \item Theorem~\ref{thm:adaptive}: median estimate $\shat_t$ adds
    $O(\sqrt{\log T/W})$ overhead relative to the oracle scale.
  \item Theorem~\ref{thm:coupled}: composing $\varphi(\phat_t)$ with $\shat_t$
    \textup{(}Torch: $\varphi(q)=q/2$\textup{)} keeps the same exponent;
    constants absorb $\varphi(p)$.
  \item Growing $\tau_t$ \textup{(}Theorem~\ref{thm:growing}\textup{)} trades a
    weaker exponent $(p-1)/p$ for not tuning $\tau$ from known $T$ in the same
    way; see the theorem for the $\eta_0(T)$ caveat.
  \item Clamping $\phat_t$ caps worst-case threshold drift when Hill is off-model.
\end{enumerate}


% ===========================================================================
\appendix
\section{High-Probability Convergence Bound}
\label{sec:highprob}
% ===========================================================================

Appendix~\ref{sec:highprob} records a high-probability analogue of
Theorem~\ref{thm:fixed} via Freedman's inequality.

\begin{theorem}[High-probability fixed-schedule convergence]\label{thm:highprob}
Under the same conditions as Theorem~\ref{thm:fixed}, for any $\delta \in (0,1)$:
with probability at least $1 - \delta$,
\begin{equation}\label{eq:highprob_rate}
  \frac{1}{T}\sum_{t=1}^T \norm{\nabla f(x_t)}^2
  \leq C_{\mathrm{rate}} \cdot T^{-\frac{2(p-1)}{3p-2}}
  + C_{\mathrm{hp}} \cdot T^{-\frac{2(p-1)}{3p-2}} \cdot \sqrt{\log(1/\delta)},
\end{equation}
where $C_{\mathrm{rate}}$ is from Theorem~\ref{thm:fixed} and $C_{\mathrm{hp}}$
depends on $L, \Delta_f, \sigma, M, p$ but not on $T$ or $\delta$.
\end{theorem}

\begin{proof}
\textbf{Step 1: Martingale decomposition.}
Define the martingale difference sequence:
\[
  M_t := f(x_{t+1}) - \E[f(x_{t+1}) \mid \Filt_t], \qquad t = 1, \ldots, T.
\]
By construction, $\E[M_t \mid \Filt_t] = 0$. From the descent
lemma~\eqref{eq:descent_clip}:
\[
  f(x_{t+1}) = f(x_t) - \eta\inner{\nabla f(x_t)}{g_t^c}
    + \frac{L\eta^2}{2}\norm{g_t^c}^2,
\]
so:
\begin{align}
  M_t &= -\eta\inner{\nabla f(x_t)}{g_t^c - \E[g_t^c \mid \Filt_t]}
    + \frac{L\eta^2}{2}\!\left(\norm{g_t^c}^2
      - \E[\norm{g_t^c}^2 \mid \Filt_t]\right) \notag\\
  &= -\eta\inner{\nabla f(x_t)}{\xi_t}
    + \frac{L\eta^2}{2}\!\left(\norm{g_t^c}^2 - \E[\norm{g_t^c}^2 \mid \Filt_t]\right),
    \label{eq:M_decomp}
\end{align}
where $\xi_t := g_t^c - \E[g_t^c \mid \Filt_t]$ is the zero-mean residual.

\textbf{Step 2: Bounded increments.}
Since $\norm{g_t^c} \leq \tau$ and $\norm{\nabla f(x_t)} \leq M$:
\begin{itemize}[noitemsep]
  \item $\abs{\inner{\nabla f(x_t)}{\xi_t}} \leq M\norm{\xi_t}
    \leq M \cdot 2\tau = 2M\tau$.
  \item $\abs{\norm{g_t^c}^2 - \E[\norm{g_t^c}^2 \mid \Filt_t]} \leq \tau^2$.
\end{itemize}
Therefore:
\begin{equation}\label{eq:M_bound}
  \abs{M_t} \leq 2\eta M\tau + \frac{L\eta^2\tau^2}{2}
  = \eta\tau\!\left(2M + \frac{L\eta\tau}{2}\right) =: c_M.
\end{equation}

\textbf{Step 3: Conditional variance bound.}
The conditional variance of $M_t$ is:
\[
  \E[M_t^2 \mid \Filt_t]
  \leq 2\eta^2 M^2 \E[\norm{\xi_t}^2 \mid \Filt_t]
    + \frac{L^2\eta^4}{2}\E\!\left[\!\left(\norm{g_t^c}^2
      - \E[\norm{g_t^c}^2]\right)^{\!2} \mid \Filt_t\right],
\]
using $(a + b)^2 \leq 2a^2 + 2b^2$.
Since $\E[\norm{\xi_t}^2] \leq \E[\norm{g_t^c}^2] \leq V(\tau)^2$ and
$\E[(\norm{g_t^c}^2 - \E[\norm{g_t^c}^2])^2] \leq \E[\norm{g_t^c}^4]
\leq \tau^2 V(\tau)^2$:
\begin{equation}\label{eq:cond_var}
  \E[M_t^2 \mid \Filt_t]
  \leq 2\eta^2 M^2 V^2 + \frac{L^2\eta^4\tau^2 V^2}{2}
  = \eta^2 V^2\!\left(2M^2 + \frac{L^2\eta^2\tau^2}{2}\right)
  =: \sigma_M^2.
\end{equation}

The \emph{predictable quadratic variation} is:
\[
  W_T := \sum_{t=1}^T \E[M_t^2 \mid \Filt_t] \leq T\sigma_M^2.
\]

\textbf{Step 4: Apply Freedman's inequality.}
Freedman's inequality (see Freedman, 1975; or Tropp, 2011, Theorem~1.6):
For a martingale $S_T = \sum_{t=1}^T M_t$ with $\abs{M_t} \leq c_M$ and
predictable quadratic variation $W_T$, for any $s, w > 0$:
\[
  \Prob\!\left(S_T \geq s \text{ and } W_T \leq w\right)
  \leq \exp\!\left(-\frac{s^2}{2(w + c_M s/3)}\right).
\]
Setting $w = T\sigma_M^2$ and solving for $s$ such that the RHS $= \delta/2$:
\[
  s = \sqrt{2T\sigma_M^2 \ln(2/\delta)} + \frac{2c_M \ln(2/\delta)}{3}.
\]
A symmetric argument applies for $-S_T$. By a union bound: with probability
$\geq 1 - \delta$,
\begin{equation}\label{eq:freedman_applied}
  \abs{S_T} \leq \sqrt{2T\sigma_M^2 \ln(2/\delta)} + \frac{2c_M \ln(2/\delta)}{3}.
\end{equation}

\textbf{Step 5: Convert to convergence bound.}
From the descent telescope (summing~\eqref{eq:descent_final}):
\[
  f(x_{T+1}) \leq f(x_1) - \frac{\eta}{2}\sum_{t=1}^T \norm{\nabla f(x_t)}^2
    + \frac{T\eta}{2}B^2 + \frac{TL\eta^2}{2}V^2 + S_T.
\]
Since $f(x_{T+1}) \geq f^*$:
\[
  \frac{\eta}{2}\sum_{t=1}^T \norm{\nabla f(x_t)}^2
  \leq \Delta_f + \frac{T\eta}{2}B^2 + \frac{TL\eta^2}{2}V^2 + S_T.
\]
Dividing by $T\eta/2$ and using~\eqref{eq:freedman_applied}:
\begin{equation}\label{eq:hp_full}
  \frac{1}{T}\sum_{t=1}^T \norm{\nabla f(x_t)}^2
  \leq \underbrace{\frac{2\Delta_f}{\eta T} + B^2 + L\eta V^2}_{\text{in-expectation bound}}
  + \underbrace{\frac{2\sqrt{2T\sigma_M^2 \ln(2/\delta)}}{T\eta}
    + \frac{4c_M\ln(2/\delta)}{3T\eta}}_{\text{high-probability correction}}.
\end{equation}

\textbf{Step 6: Evaluate the correction terms at optimal $\tau^*$, $\eta^*$.}

\emph{First correction term.}
\[
  \frac{2\sigma_M\sqrt{2T\ln(2/\delta)}}{T\eta}
  = \frac{2\sigma_M\sqrt{2\ln(2/\delta)}}{\eta\sqrt{T}}.
\]
Recall $\sigma_M = \eta V\sqrt{2M^2 + L^2\eta^2\tau^2/2}$. With
$\eta^* = \Theta(T^{-p/(3p-2)})$ and $\tau^* = \Theta(T^{1/(3p-2)})$:
\[
  \frac{\sigma_M}{\eta\sqrt{T}} = \frac{V\sqrt{2M^2 + O(L^2\eta^2\tau^2)}}{\sqrt{T}}.
\]
The dominant part is $MV/\sqrt{T}$. With $V = V(\tau^*)
= \Theta(\tau^{*(2-p)/2}) = \Theta(T^{(2-p)/(2(3p-2))})$:
\[
  \frac{MV}{\sqrt{T}} = \Theta\!\left(T^{(2-p)/(2(3p-2)) - 1/2}\right)
  = \Theta\!\left(T^{(2-p-3p+2)/(2(3p-2))}\right)
  = \Theta\!\left(T^{-(2p-2)/(3p-2)}\right)
  = \Theta\!\left(T^{-2(p-1)/(3p-2)}\right).
\]
The first correction matches the in-expectation rate!

\emph{Second correction term.}
\[
  \frac{4c_M\ln(2/\delta)}{3T\eta}
  = \frac{4\eta\tau(2M + O(L\eta\tau))\ln(2/\delta)}{3T\eta}
  = \frac{4\tau(2M + O(L\eta\tau))\ln(2/\delta)}{3T}.
\]
With $\tau = \Theta(T^{1/(3p-2)})$:
$\tau/T = \Theta(T^{1/(3p-2) - 1}) = \Theta(T^{-(3p-3)/(3p-2)})$.
Since $(3p-3)/(3p-2) > 2(p-1)/(3p-2)$ iff $3p - 3 > 2p - 2$ iff $p > 1$
(always true), this term decays faster than the rate. So it is dominated.

\textbf{Conclusion.}
Collecting the in-expectation bound and the first correction:
\[
  \frac{1}{T}\sum_{t=1}^T \norm{\nabla f(x_t)}^2
  \leq \left(C_{\mathrm{rate}} + C_{\mathrm{hp}}\sqrt{\log(2/\delta)}\right)
    \cdot T^{-2(p-1)/(3p-2)}
  + o\!\left(T^{-2(p-1)/(3p-2)}\right),
\]
with probability $\geq 1 - \delta$, where:
\begin{align*}
  C_{\mathrm{hp}} &= 4M\sqrt{A_V} \cdot
    \left(\frac{A_B}{2\sqrt{2L\Delta_f A_V}}\right)^{\frac{2-p}{2(3p-2)}}
    \bigg/ \sqrt{2\Delta_f / (LA_V)}.
\end{align*}
(An explicit but complicated constant depending only on problem parameters.)
\end{proof}

\begin{remark}[Comparison with Azuma--Hoeffding]
Using Azuma--Hoeffding instead of Freedman gives a correction term of
$O(\tau\sqrt{\log(1/\delta)/T})$, which at optimal $\tau^* = T^{1/(3p-2)}$
becomes $O(T^{1/(3p-2) - 1/2}\sqrt{\log}) = O(T^{(4-3p)/(2(3p-2))}\sqrt{\log})$.
For $p < 4/3$, this \emph{grows} with $T$, making Azuma--Hoeffding useless.
Freedman's inequality exploits the conditional variance (which is $O(\eta^2 V^2)$,
much smaller than the worst-case increment $c_M^2$), yielding a correction
that matches the in-expectation rate for all $p \in (1, 2]$.
\end{remark}


% ===========================================================================
\begin{thebibliography}{9}

\bibitem{zhang2020}
  J.\ Zhang, T.\ He, S.\ Sra, A.\ Jadbabaie.
  \textit{Why Gradient Clipping Accelerates Training: A Theoretical Justification
  for Adaptivity}. ICLR 2020.

\bibitem{mason1982}
  D.M.\ Mason.
  \textit{Laws of large numbers for sums of extreme values}.
  Annals of Probability, 10(3):754--764, 1982.

\bibitem{haeusler1985}
  E.\ Haeusler and J.L.\ Teugels.
  \textit{On asymptotic normality of Hill's estimator for the exponent of
  regular variation}. Annals of Statistics, 13(2):743--756, 1985.

\bibitem{vershynin2018}
  R.\ Vershynin.
  \textit{High-Dimensional Probability: An Introduction with Applications in
  Data Science}. Cambridge University Press, 2018.

\bibitem{beirlant2004}
  J.\ Beirlant, Y.\ Goegebeur, J.\ Segers, J.\ Teugels.
  \textit{Statistics of Extremes: Theory and Applications}. Wiley, 2004.

\bibitem{cutkosky2021}
  A.\ Cutkosky and H.\ Mehta.
  \textit{High-dimensional robust mean estimation via gradient descent}.
  ICML 2021.

\end{thebibliography}

\end{document}
